\chapter{Экспериментальное приложение}
\label{sec:supplement}

Для каждой главы здесь собраны её главные утверждения и то, на чём они держатся: опыт, в котором это измерено, что именно измерено и где это опубликовано. Статус в правом столбце говорит, насколько утверждение твёрдо: \emph{измерено} --- есть прямой опыт; \emph{теория} --- математическое следствие, выведенное в главе или в классической работе; \emph{модель} --- результат расчёта по модели книги (скрипты в папке \texttt{models/}); \emph{оценка} --- число порядка величины без прямого измерения; \emph{гипотеза} --- правдоподобное объяснение, которое опытом ещё не доказано. Где опыта нет, так и сказано.

\section{Глава~\ref{sec:neurontypes}. Типы нейронов}

Числа главы держатся на прямых подсчётах клеток в мозге человека там, где такие подсчёты есть, правило «один нейрон --- один медиатор» --- на клеточных атласах и опытах, которые нашли и исключения, роль \nt{ДОФА} --- на записях импульсов, а роли остальных широковещательных каналов --- на гипотезах.

{\footnotesize
\setlength{\tabcolsep}{3pt}\renewcommand{\arraystretch}{1.15}
\begin{longtable}{>{\raggedright}p{0.5cm}>{\raggedright}p{4.0cm}>{\raggedright}p{5.6cm}>{\raggedright}p{2.3cm}>{\raggedright\arraybackslash}p{1.7cm}}
\toprule
№ & Утверждение & Опыт и что измерено & Источник & Статус \\
\midrule\endfirsthead
\toprule
№ & Утверждение & Опыт и что измерено & Источник & Статус \\
\midrule\endhead
1 & В мозге человека около $8.6\cdot10^{10}$ нейронов, и почти все \nt{ГЛУТ}-нейроны --- мелкие нейроны мозжечка (табл.~\ref{tab:types}) & Изотропный фракционатор: ткань растворяют в однородную взвесь ядер и считают ядра с нейронным маркером NeuN. У взрослых мужчин $86.1\pm8.1$ млрд нейронов; в коре больших полушарий из них лишь $19\,\%$, основная масса --- в мозжечке & \cite{Azevedo2009} & измерено \\
2 & У почти каждого нейрона один главный медиатор (утверждение~\ref{prop:dale}), но есть исключения & Транскриптомный атлас мозга мыши: около $4\cdot10^6$ клеток, $5322$ кластера; кластерам приписан медиатор по генам синтеза и упаковки. Исключения измерены прямо: \nt{ДОФА}-нейроны в срезах выбрасывают ещё и ГАМК через тот же везикулярный транспортёр и тормозят нейроны стриатума; у части \nt{ДОФА}-аксонов глутамат и дофамин выходят из разных окончаний одного провода (электронная микроскопия и оптогенетика) & \cite{Strata1999,Yao2023,Tritsch2012,Zhang2015} & измерено \\
3 & Весь столбец матрицы весов одного знака~\eqref{eq:dalesign} & Вывод в главе из утверждения~\ref{prop:dale} и того, что рецепторы глутамата почти все возбуждающие, а ГАМК --- тормозные. Прямой проверки «все получатели одного нейрона получают вес одного знака» как общего правила нет; контрпримеры --- совместный выброс ГАМК \nt{ДОФА}-нейронами (строка~2) и получатели с рецепторами разного знака (строка~11) & вывод в главе & теория \\
4 & От трёх четвертей до четырёх пятых нейронов коры --- \nt{ГЛУТ}, от пятой части до четверти --- \nt{ГАМК} & Иммуноокраска на ГАМК в столбиках шириной $50$\,мкм через всю толщу коры в $10$ областях коры обезьяны: в $8$ областях ГАМК-нейроны составляют около $25\,\%$ всех нейронов, в первичной зрительной коре --- меньше $20\,\%$ & \cite{Hendry1987} & измерено \\
5 & У \nt{ГЛУТ}-нейрона около $10^4$ выходов & Стереология на вскрытии: в новой коре молодых мужчин $1.64\cdot10^{14}$ синапсов (электронная микроскопия, дисектор) и около $2.3\cdot10^{10}$ нейронов. Отношение даёт $\sim7\cdot10^3$ синапсов на нейрон; в среднем число входов равно числу выходов & \cite{Tang2001synapses,Pakkenberg1997} & измерено \\
6 & Выход \nt{ДОФА}-нейрона ветвится на $10^5$--$10^6$ окончаний; это оценка по охвату мишени, а не подсчёт & Вирусная метка мембраны отдельных \nt{ДОФА}-нейронов крысы и полная реконструкция аксона: один нейрон покрывает от $0.45$ до $5.7\,\%$ (в среднем $2.7\,\%$) объёма стриатума. Измерен охват, а не число окончаний: оно выведено из охвата по порядку величины, прямого счёта окончаний нет. Для \nt{СЕРО} и \nt{НОРА} числа окончаний --- грубые оценки & \cite{Matsuda2009} & измерено \\
7 & Широковещательных нейронов мало: \nt{ДОФА} $\sim5\cdot10^5$ (главная из двух групп, нижняя граница), \nt{СЕРО} $\sim3\cdot10^5$ (сумма двух частичных подсчётов), \nt{ГИСТ} $\sim6\cdot10^4$ (табл.~\ref{tab:types}, рис.~\ref{fig:typepie}) & Стереологические подсчёты у человека: $550\,000$ пигментированных нейронов в чёрной субстанции (без вентральной покрышки); $235\,000$ нейронов в дорсальном ядре шва (все нейроны ядра) и ещё $125\,000$ серотониновых нейронов в покрышке моста; около $64\,000$ гистаминовых нейронов гипоталамуса. Для \nt{НОРА}, \nt{ГЛИЦ}, \nt{АЦЕТ} и \nt{АДРЕ} прямого подсчёта нет: в таблице это грубые оценки по порядку величины & \cite{Pakkenberg1991,Baker1990,Baker1991,Airaksinen1991} & измерено \\
8 & Глутамат в щели взлетает примерно до миллимоля и спадает за $\sim1\ms$ & По вытеснению быстро отсоединяющегося конкурентного блокатора с рецепторов NMDA во время синаптической передачи (культура нейронов гиппокампа): пик $1.1$\,мМ, постоянная спада $1.2\ms$ & \cite{Clements1992} & измерено \\
9 & В работающей коре возбуждающий и тормозящий токи почти уравновешены, нейрон сидит у порога & Теория: сеть с сильными связями сама приходит в уравновешенный режим. Опыт: в префронтальной коре хорька in vivo во время «вверх»-состояний потенциал реверсии синаптического тока держится около $-37$\,мВ сотни миллисекунд, хотя проводимость меняется на $\sim21$\,нСм: возбуждение и торможение растут и падают вместе & \cite{vanVreeswijk1996,Haider2006} & измерено \\
10 & \nt{ДОФА}-нейроны сообщают ошибку предсказания награды~\eqref{eq:rpe} (рис.~\ref{fig:rpe}) & Импульсы \nt{ДОФА}-нейронов обезьяны: пачка на неожиданный сок; после обучения --- пачка на сигнал и нет ответа на предсказанный сок; провал частоты, когда обещанный сок не дан. В количественном опыте частота пропорциональна разности текущей награды и взвешенного среднего прошлых, но только для положительных ошибок. Само уравнение~\eqref{eq:rpe} --- алгоритм временных разностей & \cite{Schultz1997,Bayer2005,SuttonBarto2018} & измерено \\
11 & Знак и скорость задаёт рецептор: ионотропные --- доли миллисекунды, метаботропные --- гораздо медленнее (утверждение~\ref{prop:receptor}) & Быстрые импульсы глутамата на кусочки мембраны нейронов гиппокампа: ток через ионотропные рецепторы нарастает (от $20$ до $80\,\%$) за $0.2$--$0.6\ms$ и спадает за $2$--$3\ms$. Медленный тормозный потенциал пирамидных нейронов снимается избирательным блокатором метаботропного рецептора ГАМК. Два семейства рецепторов дофамина с противоположным действием; в стриатуме нейронам с рецептором D1 дофамин нужен для усиления синапсов, нейронам с D2 --- для ослабления & \cite{Colquhoun1992,DutarNicoll1988,Beaulieu2011,Shen2008} & измерено \\
12 & Широковещательный канал --- не один провод, а жгут из немногих линий (раздел~\ref{sec:buses}) & Вирусно-генетическая разметка серотониновых нейронов дорсального ядра шва мыши: нейроны, посылающие выход в миндалину и в лобную кору, лежат в разных местах ядра, ветвятся в дополняющие друг друга области и противоположно отвечают на неприятный раздражитель. Обзор: подгруппы нейронов голубого пятна (\nt{НОРА}) кодируют разные процессы & \cite{Ren2018,Poe2020} & измерено \\
13 & \nt{НОРА} задаёт коэффициент усиления $g$ & Гипотеза адаптивного усиления; модель сети, в которой катехоламины повышают крутизну передаточной характеристики. Прямого измерения «$g$ растёт с норадреналином» во всей сети нет; измерено, что диаметр зрачка следует за импульсами нейронов голубого пятна обезьяны & \cite{AstonJones2005,ServanSchreiber1990,Joshi2016} & гипотеза \\
14 & \nt{АЦЕТ} переключает «запись/чтение» и задаёт скорость обучения & Гипотеза. Опора: в срезах обонятельной коры крысы агонисты ацетилхолина ($100$\,мкМ) ослабляют внутренние, «вспоминающие» синапсы больше чем на $60\,\%$, а входные --- меньше чем на $15\,\%$ & \cite{Hasselmo2006,HasselmoBower1992} & гипотеза \\
15 & \nt{СЕРО} задаёт коэффициент дисконтирования $\gamma$; серотониновые нейроны работают ровно, несколько импульсов в секунду, и почти молчат в одной из фаз сна & Гипотеза «серотонин $=\gamma$». Сильнейший довод: оптогенетическая активация серотониновых нейронов дорсального ядра шва мыши уменьшает число преждевременных отказов от ожидания отложенной награды. Частота импульсов и молчание во сне с быстрыми движениями глаз измерены (обзор записей) & \cite{Doya2002,Miyazaki2014,JacobsAzmitia1992} & гипотеза \\
\bottomrule
\end{longtable}}

\section{Глава~\ref{sec:glutcomputer}. Что вычисляют \nt{GLUT}-нейроны}

Логические утверждения главы --- теоремы о схемах из неотрицательных весов, выведенные в самой главе; опыт подтверждает модель нейрона и то, что схемы, собранные по этим теоремам (детектор совпадений, линия задержки, самоподдерживающаяся группа, цепочка), в мозге действительно встречаются.

{\footnotesize
\setlength{\tabcolsep}{3pt}\renewcommand{\arraystretch}{1.15}
\begin{longtable}{>{\raggedright}p{0.5cm}>{\raggedright}p{4.0cm}>{\raggedright}p{5.6cm}>{\raggedright}p{2.3cm}>{\raggedright\arraybackslash}p{1.7cm}}
\toprule
№ & Утверждение & Опыт и что измерено & Источник & Статус \\
\midrule\endfirsthead
\toprule
№ & Утверждение & Опыт и что измерено & Источник & Статус \\
\midrule\endhead
1 & Нейрон --- RC-цепь с порогом и сбросом~\eqref{eq:glutneuron} & В тело пирамидных нейронов коры крысы (срезы) подавали шумовой ток, похожий на синаптический поток в живом мозге. Зависимость средней частоты от среднего и разброса тока описывается интегрирующим нейроном с утечкой, если добавить медленную адаптацию частоты. Диапазоны $\tau$ и задержек приведены в главе без ссылки на опыт & \cite{Rauch2003} & измерено \\
2 & Модель~\eqref{eq:glutneuron} --- упрощение: ветви входов сами выдают местные импульсы & Запись прямо с дендритов пирамидных нейронов зрительной коры мыши in vivo: зрительный раздражитель вызывает местные дендритные импульсы, зависящие от рецепторов NMDA; их подавление расширяет настройку нейрона на ориентацию & \cite{Smith2013} & измерено \\
3 & Окно совпадения $\Delta_{\max}=\tau\ln\frac{w}{\theta-w}$~\eqref{eq:window}; при $\theta=1.5w$ оно равно $0.7\tau$ & Следствие модели~\eqref{eq:glutneuron}. Прямое измерение узкого окна суммирования есть для нейронов с быстрым торможением (глава~\ref{sec:gaba}, строка~3 её таблицы) & вывод в главе & теория \\
4 & Сеть из одних \nt{GLUT}-нейронов вычисляет только монотонные функции входов (утверждение~\ref{prop:monotone}) & Доказательство в главе: итерация от тишины с неубывающей правой частью. Это утверждение о модели; опыта, который проверял бы его на живой сети без торможения, нет & вывод в главе & теория \\
5 & Органы чувств дают пару проводов: одни клетки отвечают на усиление раздражителя, другие --- на ослабление & Сетчатка: уже на биполярных клетках сигнал колбочек расходится на включающий и выключающий каналы. Фармакологический блок включающих биполярных клеток у обезьяны: каналы идут раздельно до коры; после блока ухудшается обнаружение усиления света, но не ослабления & \cite{Schiller1992} & измерено \\
6 & С двухпроводным кодом сеть из \nt{GLUT}-нейронов вычисляет любую логическую функцию асинхронно, без общего такта, ценой вдвое большего числа нейронов (утверждение~\ref{prop:dualrail}, \eqref{eq:dualrail}, рис.~\ref{fig:dualxor}) & Двухпроводный код и определение готовности по самому сигналу --- стандартная техника асинхронных схем, нечувствительных к задержкам. Схема исключающего ИЛИ из шести нейронов собрана в главе. Опыта, показывающего, что мозг вычисляет логические функции именно двухпроводным кодом, нет & \cite{Sparso2001} & теория \\
7 & Наибольшая разница времени прихода звука в уши у человека $\approx0.6\ms$, а мозг различает около десяти микросекунд & Слушатели в опыте с двумя вариантами ответа: порог различения разницы времени (75\,\% правильных) --- $9$\,мкс для шума в полосе $150$--$1700$\,Гц, $11$\,мкс для тона $1$\,кГц, $28$\,мкс для щелчка. Предел $0.6\ms$ --- расчёт главы, $0.22\,\text{м}/343\,\text{м/с}$ & \cite{KlumppEady1956} & измерено \\
8 & У птиц разница времён превращается в место через линии задержки и детекторы совпадений~\eqref{eq:itd} (рис.~\ref{fig:itd}) & Сипуха: аксоны от двух ушей входят в ядро детекторов совпадений с противоположных сторон; время прихода импульсов вдоль аксона меняется упорядоченно с глубиной, и разброс задержек через ядро ($160$\,мкс на $700$\,мкм) совпадает с диапазоном межушных задержек & \cite{Carr1990} & измерено \\
9 & У млекопитающих ряда задержек не нашли; ту же задачу решают детекторы совпадений с точно рассчитанным торможением от \nt{GLYC}-нейронов & Запись in vivo в медиальной верхней оливе песчанки: ответы не образуют карты направлений; подведение глицина и его блокатора через микропипетку показывает, что точно рассчитанное во времени \emph{глициновое} торможение задаёт рабочий диапазон задержек. Торможение здесь от \nt{GLYC}, а не от \nt{GABA} & \cite{Brand2002,Grothe2010} & измерено \\
10 & Группа с сильной обратной связью --- триггер; самоподдерживающаяся активность держится секунды после раздражителя (рис.~\ref{fig:bistable}) & Префронтальная кора обезьяны в задаче с отложенным движением глаз к запомненной точке (задержка $1$--$6$\,с): $87$ из $170$ нейронов, связанных с задачей, меняют частоту во время задержки, у $79\,\%$ из них это зависит от направления на запомненную точку. Что эту активность держит обратная связь с двумя устойчивыми состояниями, --- теория & \cite{Funahashi1989,Wang2001} & измерено \\
11 & Бит записывается, если приток длится дольше $\approx0.7\tau$ (рис.~\ref{fig:recurrentgroup}б) & Расчёт уравнения $\tau\,\dd r/\dd t=-r+f(wr+I)$ методом Эйлера при $w=1$, $I=0.6$; скрипта в \texttt{models/} нет. Опыта нет & расчёт в главе & модель \\
12 & Память можно хранить в облегчении синапсов, почти без импульсов & Теория: остаточный кальций в окончаниях держит запись около секунды. Опора: в медиальной префронтальной коре хорька (запись с нескольких нейронов сразу) есть подсеть пирамидных нейронов, связанных облегчающимися синапсами с долгим последействием. Обзор наблюдений «тихих» промежутков при удержании в памяти. Какой способ кора использует когда --- открытый вопрос & \cite{Mongillo2008,WangMarkram2006,Stokes2015} & гипотеза \\
13 & Память стирается усталостью: запас медиатора истощается & Парные записи пирамидных нейронов коры: ответ синапса падает при частых импульсах, скорость падения задаётся вероятностью выброса. Что именно это гасит самоподдерживающуюся группу, прямо не показано & \cite{Tsodyks1997} & измерено \\
14 & Цепочка групп --- таймер, запускаемый событием; у певчих птиц нейроны срабатывают строго по очереди & Запись опознанных нейронов высшего вокального центра зебровой амадины во сне и при пении: каждый нейрон, посылающий выход в моторное ядро, выдаёт одну короткую пачку в один точный момент песни, и нейроны срабатывают последовательно & \cite{Hahnloser2002} & измерено \\
\bottomrule
\end{longtable}}

\section{Глава~\ref{sec:gaba}. \nt{ГЛУТ} и \nt{ГАМК} вместе}

Логические и динамические утверждения главы выведены в ней самой из линейного анализа и закона Ома; опыт показывает, что схемы, на которые они опираются, --- запрет с задержкой, прямое торможение вслед за возбуждением, стабилизация торможением, ритм петли \nt{ГЛУТ}--\nt{ГАМК}, --- в мозге действительно работают.

{\footnotesize
\setlength{\tabcolsep}{3pt}\renewcommand{\arraystretch}{1.15}
\begin{longtable}{>{\raggedright}p{0.5cm}>{\raggedright}p{4.0cm}>{\raggedright}p{5.6cm}>{\raggedright}p{2.3cm}>{\raggedright\arraybackslash}p{1.7cm}}
\toprule
№ & Утверждение & Опыт и что измерено & Источник & Статус \\
\midrule\endfirsthead
\toprule
№ & Утверждение & Опыт и что измерено & Источник & Статус \\
\midrule\endhead
1 & \nt{ГАМК}-нейронов в коре около пятой части & Иммуноокраска на ГАМК в $10$ областях коры обезьяны: около $25\,\%$ нейронов в $8$ областях, меньше $20\,\%$ в первичной зрительной коре. Обзор разнообразия тормозящих нейронов коры & \cite{Hendry1987,Markram2004} & измерено \\
2 & Что делает торможение, во многом решает место посадки: дендрит, тело, место рождения импульса, окончание чужого провода & Обзор: классы \nt{ГАМК}-нейронов коры различаются мишенью на получателе. Одновременная запись с тела и дендрита пирамидного нейрона: прямое торможение гораздо сильнее на теле, чем на дендритах. Торможение на окончании чужого провода в спинном мозге уменьшает выброс медиатора одной входной линии (обзор измерений) & \cite{Markram2004,Pouille2001,Rudomin1999} & измерено \\
3 & Смена знака через посредника стоит одной задержки, около миллисекунды; торможение, пришедшее вслед за возбуждением, сужает окно совпадения & Пирамидные нейроны гиппокампа крысы в срезах: торможение через одного посредника приходит вслед за прямым возбуждением с короткой задержкой, и окно, в котором два возбуждающих входа суммируются до порога, меньше $2\ms$. На дендритах, где это торможение слабее, окно шире & \cite{Pouille2001} & измерено \\
4 & Запрет $a\wedge\bar b$~\eqref{eq:veto} с ИЛИ и постоянной единицей функционально полон (утверждение~\ref{prop:gabacomplete}) & Доказательство в главе. Классический результат: сеть пороговых элементов с тормозящими входами реализует любое логическое выражение. Утверждение о модели, опыта нет & \cite{McCullochPitts1943} & теория \\
5 & Запрет с задержкой даёт детектор направления движения с оптимальной скоростью $v^*=\Delta x/d$~\eqref{eq:vstar} & Сетчатка кролика: избирательность к направлению создаётся торможением, которое при движении в «нулевом» направлении гасит возбуждение. Раздельная запись возбуждающего и тормозящего входов избирательных клеток показала, что звёздчатые амакриновые клетки (\nt{ГАМК}) тормозят их несимметрично --- только отростками, направленными в «нулевую» сторону. Формула $v^*$ --- вывод главы & \cite{Barlow1965,Fried2002} & измерено \\
6 & Шунтирующее торможение делит напряжение, а не вычитает~\eqref{eq:shunt} (рис.~\ref{fig:shunt}) & Закон Ома для делителя из трёх проводимостей при равновесном потенциале хлора около покоя; для нейрона без импульсов & вывод в главе & теория \\
7 & На частоту импульсов шунт действует вычитанием, а деление получается от торможения вместе с уравновешенным фоновым шумом & Модель: у срабатывающего нейрона ток через шунт почти постоянен, и эффект на частоту вычитающий. Опыт: в пирамидные нейроны соматосенсорной коры крысы (срезы) через динамический фиксатор вводили фоновые возбуждающие и тормозящие проводимости; одновременный уравновешенный рост обеих делит крутизну ответа без заметного сдвига частоты. Обзор: деление на общую активность встречается во многих отделах мозга & \cite{HoltKoch1997,Chance2002,Carandini2012} & измерено \\
8 & При $\beta>1$ взаимное торможение выбирает одного победителя (утверждение~\ref{prop:wta}, \eqref{eq:wta}) & Собственные числа $-(1\pm\beta)/\tau$ --- вывод в главе. Частоты $0.73$ и $0.53$ при $\beta=0.5$ --- неподвижная точка $r_{1,2}=(I_{1,2}-\beta I_{2,1})/(1-\beta^2)$; скрипта в \texttt{models/} нет. Прямого опыта в главе не приводится & вывод в главе & теория \\
9 & Сброс памяти сигналом через \nt{ГАМК}; две группы с взаимным торможением --- защёлка & Следует из рис.~\ref{fig:bistable} и утверждения~\ref{prop:wta}. Опыта нет & вывод в главе & теория \\
10 & Равновесие петли \nt{ГЛУТ}--\nt{ГАМК} устойчиво, если торможение достаточно сильное и достаточно быстрое (утверждение~\ref{prop:ei}) & Модель двух популяций~\eqref{eq:ei}; условия (i) и (ii) --- определитель и след её матрицы, выведены в главе & \cite{WilsonCowan1972} & теория \\
11 & При $w_{EE}=1.5$ и $w_{EI}w_{IE}=2$ быстрое торможение ($\tau_I=2\ms$) держит $E^*=0.67h$, медленное ($30\ms$) раскачивает колебания (рис.~\ref{fig:ei}) & Численное решение~\eqref{eq:ei}, скрипт \texttt{figures/make\_ei.py} & расчёт & модель \\
12 & Кора работает в режиме стабилизации торможением; подпись --- подтолкнёшь \nt{ГАМК}-нейроны, и их частота упадёт~\eqref{eq:isn} & Теория предсказала парадоксальный ответ. Опыт: внутриклеточные записи в первичной зрительной коре кошки --- при подавлении ответа раздражителем вокруг тормозящая проводимость не растёт, а падает вместе с возбуждающей. Оптогенетическое возбуждение PV-нейронов мыши в зрительной, соматосенсорной и моторной коре снижает частоту тормозящих нейронов, в том числе без раздражителя и при лёгком наркозе. Коэффициент $-1/3$ при числах рис.~\ref{fig:ei} --- вывод главы & \cite{Tsodyks1997isn,Ozeki2009,Sanzeni2020} & измерено \\
13 & Петля «\nt{ГЛУТ} возбуждает \nt{ГАМК}, \nt{ГАМК} тормозит \nt{ГЛУТ}» порождает быстрый ритм & Оптогенетическая стимуляция быстрых \nt{ГАМК}-нейронов в соматосенсорной коре мыши на частотах $8$--$200$\,Гц избирательно усиливает гамма-ритм ($20$--$80$\,Гц, период $12$--$50\ms$); стимуляция \nt{ГЛУТ}-нейронов усиливает только более медленные ритмы. В главе период назван $10$--$30\ms$, то есть верхняя часть этой полосы & \cite{Cardin2009,Buzsaki2012} & измерено \\
\bottomrule
\end{longtable}}

\section{Глава~\ref{sec:code}. Азбука Морзе}

Предельные оценки главы --- теория информации и счёт; измерено, где в канале возникает шум, как быстро работает мозг и сколько стоит импульс, а вопрос о том, каким кодом кора пользуется на самом деле, остаётся открытым.

{\footnotesize
\setlength{\tabcolsep}{3pt}\renewcommand{\arraystretch}{1.15}
\begin{longtable}{>{\raggedright}p{0.5cm}>{\raggedright}p{4.0cm}>{\raggedright}p{5.6cm}>{\raggedright}p{2.3cm}>{\raggedright\arraybackslash}p{1.7cm}}
\toprule
№ & Утверждение & Опыт и что измерено & Источник & Статус \\
\midrule\endfirsthead
\toprule
№ & Утверждение & Опыт и что измерено & Источник & Статус \\
\midrule\endhead
1 & Частоты \nt{ГЛУТ}-нейронов коры --- от долей до десятков импульсов в секунду, и большую часть времени нейроны работают у нижней границы & Запись через прилегающую к клетке пипетку (выбор нейрона не зависит от его активности) в слуховой коре бодрствующей крысы: в каждый момент частоту выше $20$ импульсов в секунду дают меньше $5\,\%$ нейронов; у части нейронов фоновая частота ниже $0.01$ в секунду; распределение частот логнормальное & \cite{Hromadka2008} & измерено \\
2 & Ёмкость импульсного канала $R_{\max}\approx r\log_2\frac{e}{r\Delta t}$~\eqref{eq:spikeentropy}, почти вся она --- во времени импульсов (утверждение~\ref{prop:capacity}) & Энтропия двоичной строки из клеток длины $\Delta t$. Числа главы: $8$ бит на импульс при $r=10$ в секунду и $\Delta t=1\ms$, около $5$ бит при $\Delta t=10\ms$, меньше $2$ бит в секунду у счётчика импульсов & \cite{MacKay1952} & теория \\
3 & Чувствительные нейроны передают от одного до нескольких бит на импульс, в лучших случаях до половины предела & Восстановление раздражителя по импульсам чувствительных нейронов, для которых раздражитель известен; сводка измерений в книге & \cite{Rieke1997} & измерено \\
4 & Генератор импульсов точен; точное время задают быстрые перепады входа & Нейроны коры крысы в срезах: на повторяемый ток с быстрыми колебаниями импульсы повторяются с точностью меньше $1\ms$, на постоянный ток моменты расползаются & \cite{Mainen1995} & измерено \\
5 & Синапс ненадёжен: больше половины импульсов не доходят до получателя из-за случайности выброса & Минимальная стимуляция входов пирамидных нейронов гиппокампа (срезы): больше половины импульсов не вызывают ответа. Исключены колебания порога стимуляции, сбои проведения на ветвлениях аксона и низкая температура опыта; остаётся вероятностный выброс & \cite{AllenStevens1994} & измерено \\
6 & Поток импульсов в живой коре выглядит случайным: интервалы разбросаны как у пуассоновского потока, дисперсия числа импульсов близка к среднему; часть «шума» --- сообщение о движениях животного & Записи в зрительных областях V1 и MT бодрствующей обезьяны: коэффициент вариации интервалов около $1$, отношение дисперсии числа импульсов к среднему --- как у случайного процесса; модель, суммирующая независимые малые входы, даёт разброс гораздо меньше. В зрительной коре мыши значительная часть разброса предсказывается по видеозаписи движений самого животного & \cite{Softky1993,Stringer2019} & измерено \\
7 & Частотный код медленен: ошибка $1/\sqrt{rT}$~\eqref{eq:rateerror}, а мозг узнаёт картинку за $\sim150\ms$ & Формула --- пуассоновская статистика, вывод главы. Опыт: люди решали, есть ли животное на незнакомой фотографии, показанной на $20\ms$; потенциалы мозга на «да» и «нет» расходятся примерно через $150\ms$. Разбиение этого времени на десяток ступеней по $10$--$20\ms$ --- оценка главы, не измерение & \cite{Thorpe1996} & измерено \\
8 & Усреднение по группе ограничено корреляцией: ошибка падает не больше чем в $1/\sqrt c$ раз~\eqref{eq:correlated} (утверждение~\ref{prop:pooling}) & Пары нейронов области MT обезьяны, записанные одновременно: коэффициент корреляции чисел импульсов в среднем $0.12$, и теоретический анализ показывает, что даже такая корреляция ограничивает выигрыш от группы. Теория: предел ставит только та часть корреляции, что похожа на сигнал. Модель противоположной позиции: частота группы из $50$--$100$ нейронов читается за один межимпульсный интервал, $10$--$50\ms$, и время отдельных импульсов в коре не используется & \cite{Zohary1994,MorenoBote2014,ShadlenNewsome1998} & измерено \\
9 & Число можно писать во время первого импульса~\eqref{eq:latency}, в порядке импульсов группы или в фазе местного ритма & Сетчатка саламандры: пространственная структура коротко показанной картинки записана в относительных задержках первых импульсов выходных нейронов, почти независимо от контраста. Клетки места гиппокампа крысы: при проходе через «своё» место импульсы сдвигаются всё раньше в цикле тета-ритма ($7$--$12$\,Гц), сдвиг от $100^\circ$ до $355^\circ$. Насколько кора пользуется временем импульсов --- открытый вопрос (строка~8) & \cite{Gollisch2008,OKeefe1993} & измерено \\
10 & Какой код прочитан, решает постоянная времени получателя~\eqref{eq:reader} (утверждение~\ref{prop:reader}); в активной коре нейроны ближе к детекторам совпадений & Формула~\eqref{eq:reader} --- вывод главы. Обзор записей in vivo: в активной коре проводимость мембраны в несколько раз выше, чем в покое, и постоянная времени соответственно короче. Что кора переключает код именно так, прямо не показано & \cite{Destexhe2003} & гипотеза \\
11 & Истощающийся синапс передаёт изменение частоты, по существу $\Delta r/r$~\eqref{eq:depression} (рис.~\ref{fig:depression}) & Парные записи пирамидных нейронов коры: скорость истощения задаётся вероятностью выброса; при медленном истощении ответ отражает частоту, при быстром --- совпадения. Модель на основе измеренного истощения в коре крысы: одинаковые относительные изменения частоты на быстрых и медленных входах дают одинаковый ответ, то есть синапс передаёт $\Delta r/r$. Числа $1.7\to2.2$, $6.7$ и $90\ms$ --- расчёт главы & \cite{Tsodyks1997,Abbott1997depression} & измерено \\
12 & Пачка --- «тире»: её вероятность потеряться $(1-p)^k$~\eqref{eq:burst} мала, облегчение уменьшает её ещё & Обзор: на ненадёжных синапсах одиночные импульсы часто теряются, а пачки передаются надёжно благодаря облегчению; пачка вызывает долговременные изменения синапсов. Что мозг читает пачку как отдельный знак, --- гипотеза & \cite{Lisman1997,AllenStevens1994} & гипотеза \\
13 & Быстрые \nt{ГАМК}-нейроны задают ритм и «пунктуацию»; ещё один класс тормозит тормозящих и открывает ворота (утверждение~\ref{prop:punctuation}) & Обзор свойств классов \nt{ГАМК}-нейронов коры. Оптогенетика: ритмическое возбуждение быстрых \nt{ГАМК}-нейронов вызывает гамма-ритм, и ответ на раздражитель зависит от фазы цикла. В зрительной коре мыши PV-нейроны тормозят друг друга, SST-нейроны --- все остальные классы, VIP-нейроны --- преимущественно SST. Возбуждение VIP-нейронов у бодрствующей мыши снимает торможение с \nt{ГЛУТ}-нейронов; их включает сигнал подкрепления. Разделение ролей как общее правило --- гипотеза & \cite{Tremblay2016,Cardin2009,Pfeffer2013,Pi2013} & гипотеза \\
14 & Энергия ограничивает среднюю частоту величиной порядка одного импульса в секунду~\eqref{eq:energy} & Энергетический бюджет серого вещества грызуна по анатомическим и физиологическим данным: импульсы и синаптические токи --- $47\,\%$ и $34\,\%$ энергии сигнализации; одновременно могут быть активны не больше $\sim15\,\%$ нейронов. Для коры человека: не больше $\sim1\,\%$ нейронов могут быть сильно активны одновременно. Число $\sim10^9$ АТФ на импульс и мощность $\sim10$\,Вт на сигнализацию --- оценка главы на основе этих расчётов & \cite{Attwell2001,Lennie2003} & теория \\
15 & Код разрежен и подстроен под наибольшее число бит на джоуль; кора меняет точность на энергию (утверждение~\ref{prop:economy}) & Гипотеза экономного кода. Опора: у мышей при ограничении пищи в зрительной коре падает проводимость AMPA-рецепторов и синаптический расход АТФ ($-29\,\%$), частота импульсов сохраняется, а настройка на ориентацию расширяется на $32\,\%$ & \cite{Barlow1961,Padamsey2022} & гипотеза \\
\bottomrule
\end{longtable}}

\section{Глава~\ref{sec:serocomputer}. \nt{СЕРО}-нейроны}

Свойства самих \nt{СЕРО}-нейронов (ритм, автоторможение, знак по рецептору, подсистемы) измерены прямо; вычисления главы --- регулятор, фильтр, логика --- следуют из её уравнений; $\tau_c$, коэффициент диффузии серотонина и число мест выброса --- оценки; роль петли как регулятора уровня в мозге --- гипотеза (в ядре шва автоторможение точечное и даёт лишь короткие паузы), а роль \nt{СЕРО} как горизонта --- гипотеза с поведенческими опытами в её пользу.

{\footnotesize
\setlength{\tabcolsep}{3pt}\renewcommand{\arraystretch}{1.15}
\begin{longtable}{>{\raggedright}p{0.5cm}>{\raggedright}p{4.0cm}>{\raggedright}p{5.6cm}>{\raggedright}p{2.3cm}>{\raggedright\arraybackslash}p{1.7cm}}
\toprule
№ & Утверждение & Опыт и что измерено & Источник & Статус \\
\midrule\endfirsthead
\toprule
№ & Утверждение & Опыт и что измерено & Источник & Статус \\
\midrule\endhead
1 & Знак действия серотонина выбирает рецептор получателя (утверждение~\ref{prop:receptor}) & Рецепторы серотонина разделены на семейства по фармакологии и механизму: 5-HT$_{1A}$ через G-белок открывает калиевые каналы и тормозит клетку, 5-HT$_{2A}$ и ионотропный 5-HT$_3$ возбуждают. Один медиатор даёт противоположные ответы в разных клетках. & \cite{BarnesSharp1999} & измерено \\
2 & \nt{СЕРО}-нейрон в бодрствовании ровно выдаёт несколько импульсов в секунду & Запись нейронов дорсального ядра шва у свободно движущихся кошек: медленный ритмичный разряд $2.8\pm0.2$ имп/с в спокойном бодрствовании; в медленном сне частота падает на $34$--$68\,\%$, в быстром --- на $98\,\%$. Под наркозом у крыс --- $1.7\pm0.5$ имп/с, очень регулярно. & \cite{TrulsonJacobs1979, GagnonParent2014, JacobsAzmitia1992} & измерено \\
3 & \nt{СЕРО}-нейрон --- ритмоводитель: толчок на каждый импульс не нужен, но нужен ровный фоновый приток (в срезе --- норадреналин через $\alpha_1$-рецепторы, в организме --- от \nt{НОРА}) & В срезе мозга клетки разряжаются медленно и ровно, после каждого импульса --- следовая гиперполяризация на $200$--$800\ms$. Спонтанно работающих клеток в срезе меньше, чем в организме: ровному ритму нужен тонический вход норадреналина через $\alpha_1$-рецепторы, блокада которых его гасит. & \cite{VandermaelenAghajanian1983} & измерено \\
4 & Почти все рецепторы метаботропные; ответ медленный: нарастает и спадает в пределах примерно секунды & Все семейства рецепторов, кроме 5-HT$_3$, связаны с G-белками. В срезе вызванный выброс серотонина даёт через 5-HT$_{1A}$ тормозящий ток, который нарастает и спадает в пределах секунды. & \cite{BarnesSharp1999, CourtneyFord2016} & измерено \\
5 & В областях назначения серотонин выходит в объём, а не только в щель; в самом ядре шва торможение соседей идёт точка-точка & Быстрая циклическая вольтамперометрия в срезах: выброс на один импульс одинаков в области с синапсами и в ядре шва, где синапсов почти нет; не зависит от блокады захвата и рецепторов; молекул на окончание намного меньше, чем переносчиков и рецепторов, а внеклеточная концентрация совпадает со сродством главного рецептора. В ядре шва торможение через 5-HT$_{1A}$ идёт точка-точка: переносчик не даёт серотонину накопиться вне щели. & \cite{Bunin1998, CourtneyFord2016} & измерено \\
6 & Концентрация по~\eqref{eq:seroneuron} убывает из-за откачки; $\tau_c$ порядка секунды --- оценка & Вольтамперометрия в живом мозге: внеклеточный серотонин ограничивают прежде всего обратный захват и разрушение, а не синтез. Прямого измерения $\tau_c$ в цитируемых работах нет; порядок величины --- оценка главы. & \cite{Hashemi2012serotonin} & измерено; $\tau_c$ --- оценка \\
7 & НЕ и ИЛИ-НЕ на одном ритмоводителе; полнота \nt{СЕРО}-логики (утверждение~\ref{prop:serocomplete}, рис.~\ref{fig:serogates}) & Следует из~\eqref{eq:seroneuron}: тормозимый ритмоводитель --- ИЛИ-НЕ, а ИЛИ-НЕ функционально полон. Опыта, где из \nt{СЕРО}-нейронов собрана логика, нет; условие раздельных объёмов в мозге не выполняется. & вывод в главе & теория \\
8 & Серотонин тормозит сам \nt{СЕРО}-нейрон через рецепторы на нём & У наркотизированных крыс агонисты 5-HT$_{1A}$ внутривенно и ионофоретически тормозят разряд нейронов дорсального шва с той же зависимостью доза--ответ, что и сам серотонин; агонисты 5-HT$_{1B}$ почти не действуют. В срезе эндогенный серотонин соседних клеток даёт через 5-HT$_{1A}$ ток и короткую паузу в разряде, не меняя средней частоты. & \cite{Sprouse1987, CourtneyFord2016} & измерено \\
9 & В модели петля через общий объём --- регулятор уровня: разброс и постоянная времени уменьшаются в $1+G$ раз~\eqref{eq:feedback}; что петля так работает в мозге, --- гипотеза & Классическая формула усилителя с отрицательной обратной связью; в главе выведена заново. Усиление петли $G$ у \nt{СЕРО}-системы не измерено. Измерено: в срезе автоторможение идёт точка-точка, даёт короткую паузу и не меняет средней частоты. & \cite{Black1934feedback, CourtneyFord2016} & теория; в мозге --- гипотеза \\
10 & Концентрация --- скользящее среднее частоты, фильтр нижних частот~\eqref{eq:lowpass}, рис.~\ref{fig:lowpass} & Решение линейного уравнения~\eqref{eq:seroneuron}; рисунок --- расчёт по этой формуле при $\tau_c=1\,\text{с}$. Отдельной модели в \texttt{models/} нет. & вывод в главе & теория \\
11 & Извилистость щелей замедляет диффузию малой молекулы примерно в $2.6$ раза & Метод точечного источника: ионофорез тетраметиламмония и запись его концентрации ионоселективным электродом в срезах и в живом мозге. Извилистость $\lambda\approx1.6$, то есть эффективный $D$ меньше свободного в $\lambda^2\approx2.6$ раза; доля межклеточного объёма около $0.2$. Коэффициент диффузии самого серотонина в ткани в этих работах не измерялся; в главе он --- оценка. & \cite{Sykova2008} & измерено \\
12 & Длина независимого «провода» $\ell\sim\sqrt{D\tau}\approx20$~мкм~\eqref{eq:diffusionlength}; число проводов ограничено числом таких областей (утверждение~\ref{prop:volumewires}) & Закон диффузии и подстановка оценок $D$ и $\tau$ (строки~6 и~11); утверждение~\ref{prop:volumewires} --- следствие. Прямого опыта, считающего независимые каналы объёмной передачи, нет. & вывод в главе & теория \\
13 & Выход одного \nt{СЕРО}-нейрона раскинут по огромным участкам мозга; мест выброса по оценке $\sim10^5$ & Полная реконструкция отдельных нейронов дорсального шва у крысы: все восходящие аксоны сильно ветвятся, общая длина аксона одной клетки до $18.7$~см, ветви одной клетки доходят до стриатума, префронтальной коры и миндалины. Число мест выброса на клетку прямо не подсчитано. & \cite{GagnonParent2014} & измерено; $10^5$ --- оценка \\
14 & \nt{СЕРО}-нейроны делятся на несколько подсистем с разными выходами и ответами & Вирусно-генетическое мечение у мышей: нейроны, идущие к миндалине и к лобной коре, почти не пересекаются по ветвлению, получают разные входы и противоположно отвечают на неприятный стимул; включение и выключение каждой группы по-разному меняет поведение. & \cite{Ren2018} & измерено \\
15 & Условие терпения $D<\tau_\gamma\ln(R/r_0)$~\eqref{eq:patience} при экспоненциальном обесценивании; при измеренном, гиперболическом, --- $D<\tau_\gamma(R/r_0-1)$ & Следует из обесценивания $e^{-D/\tau_\gamma}$ или $1/(1+D/\tau_\gamma)$; вывод в главе. Обезьяны, выбор на задержках в секунды: обесценивание ближе к гиперболе, чем к экспоненте; ответ \nt{ДОФА}-нейронов на предвестник отложенной награды падает с задержкой тоже по гиперболе. & вывод в главе; \cite{KobayashiSchultz2008} & теория \\
16 & \nt{СЕРО} задаёт горизонт $\tau_\gamma$ (раздел~\ref{sec:horizon}) & Оптогенетическое включение \nt{СЕРО}-нейронов дорсального шва у мышей удлиняет ожидание отложенной награды и повышает упорство в поиске награды, ставшей редкой. У людей понижение серотонина (диета без триптофана) делает обесценивание отложенной награды круче. Как концентрация превращается в $\tau_\gamma$, неизвестно. & \cite{Doya2002, Miyazaki2014, Lottem2018, Schweighofer2008} & гипотеза \\
\bottomrule
\end{longtable}}

\section{Глава~\ref{sec:dofa}. Обучение с \nt{ДОФА}}

Механизмы синапса (порции, детектор совпадений, кальций, окна пластичности) и поведение дофамина как ошибки предсказания измерены прямо; то, что правила ведут к градиенту и чего стоит широковещание, --- теоремы; итог обучения выбору --- расчёт по модели книги; схема, вычисляющая ошибку, --- гипотеза.

{\footnotesize
\setlength{\tabcolsep}{3pt}\renewcommand{\arraystretch}{1.15}
\begin{longtable}{>{\raggedright}p{0.5cm}>{\raggedright}p{4.0cm}>{\raggedright}p{5.6cm}>{\raggedright}p{2.3cm}>{\raggedright\arraybackslash}p{1.7cm}}
\toprule
№ & Утверждение & Опыт и что измерено & Источник & Статус \\
\midrule\endfirsthead
\toprule
№ & Утверждение & Опыт и что измерено & Источник & Статус \\
\midrule\endhead
1 & Обратного распространения ошибки в той форме, что в искусственных сетях, в мозге не найдено & Прямого опыта нет: это утверждение об отсутствии. Обзоры разбирают, чего требует алгоритм (обратные связи, повторяющие прямые, отдельная фаза) и какие биологические приближения к нему предложены. & \cite{Lillicrap2020, WhittingtonBogacz2019} & гипотеза \\
2 & Вес раскладывается на число мест выброса, вероятность выброса и ответ на порцию: $w=N_s\,p\,A_1$~\eqref{eq:quantal} & Нервно-мышечный синапс лягушки при сниженном выбросе: ответ получателя колеблется ступенями, кратными спонтанному миниатюрному ответу на одну порцию; число порций на импульс подчиняется закону Пуассона. & \cite{delCastillo1954} & измерено \\
3 & Рецептор получателя --- детектор совпадений; кальций запускает изменение веса & Патч-клэмп нейронов мыши: ионы Mg$^{2+}$ закупоривают канал, открытый глутаматом, при потенциале покоя и уходят при деполяризации. Срезы гиппокампа: хелатор кальция в получателе блокирует долговременное усиление, а освобождение кальция светом внутри клетки само усиливает передачу. & \cite{Nowak1984, Malenka1988calcium} & измерено \\
4 & Решение исполняет любая сторона: растёт $A_1$ у получателя, меняется $p$ у отправителя & Глутамат, освобождаемый светом над одним шипиком, вызывает быстрый и стойкий рост шипика и тока через его рецепторы; усиление сопровождается встраиванием AMPA-рецепторов. Деполяризация получателя на время подавляет выброс у отправителя; эффект переносят эндоканнабиноиды, идущие назад через щель (показано на \nt{ГАМК}-входах). Кратковременное истощение зависит от вероятности выброса у отправителя. & \cite{Matsuzaki2004, Malinow2002, WilsonNicoll2001, Tsodyks1997} & измерено \\
5 & Синапс усиливается, когда отправитель и получатель активны вместе~\eqref{eq:hebb} & Кролик под наркозом: после коротких серий импульсов по входному пути гиппокампа ответ усилен на время от $30$~мин до $10$~ч. В срезе гиппокампа импульсы получателя усиливают только те синапсы, которые в тот же момент активны. & \cite{Bliss1973LTP, Kelso1986hebb} & измерено \\
6 & Правило~\eqref{eq:oja} находит главный собственный вектор корреляций входов (утверждение~\ref{prop:oja}) & Теорема; доказательство в главе. Опыта, где нейрон выучил главную компоненту своих входов, нет; механизм ограничения роста весов в синапсе не установлен. & \cite{Oja1982} & теория \\
7 & Хебб по времени: окно около $\pm20\ms$ (рис.~\ref{fig:stdp}); из повторяемых последовательностей вырастают цепочки & Пары нейронов гиппокампа в культуре: импульс получателя в пределах $20\ms$ после импульса отправителя даёт стойкое усиление, в пределах $20\ms$ до --- ослабление; оба зависят от NMDA-рецепторов. Отношение $A_-=0.6A_+$ на рисунке --- иллюстрация. Что так в сети вырастают цепочки, прямо не измерено. & \cite{BiPoo1998} & измерено \\
8 & След~\eqref{eq:threefactor}: дофамин, пришедший через $1$--$2\,\text{с}$ после совместной активности, усиливает связь, позже --- нет (утверждение~\ref{prop:threefactor}) & Срезы стриатума, раздельная оптогенетическая стимуляция глутаматных и дофаминовых входов: шипик растёт, только если дофамин приходит через $0.3$--$2\,\text{с}$ после глутаматного входа; окно задаётся быстрым подъёмом и распадом цАМФ. В коре спаривание оставляет раздельные следы для усиления и ослабления, которые моноаминовые рецепторы превращают в долговременные изменения в окне порядка секунд. & \cite{Yagishita2014, He2015eligibility} & измерено \\
9 & Окна длиной в секунды бывают и без дофамина: третий сигнал --- сильное возбуждение получателя & Живая мышь, поле CA1: одно плато-событие в получателе усиливает входы, пришедшие за секунды до и после него, и за один проход создаёт поле места. & \cite{Bittner2017} & измерено \\
10 & Средний шаг правила трёх множителей идёт по градиенту награды~\eqref{eq:perturbation} (утверждение~\ref{prop:gradient}) & Теорема о градиенте для обучения по случайным отклонениям; в главе выведена для малого шума. & \cite{Williams1992} & теория \\
11 & Шум --- это поиск: сеть учится на своих случайных отклонениях & Молодые амадины: выключение ядра, выходящего из цепи базальных ганглиев, резко и обратимо убирает изменчивость песни. Взрослые амадины: помеха подаётся на исполнения слога с высотой по одну сторону от средней, и птицы быстро смещают высоту в другую сторону --- учатся по собственному разбросу. & \cite{Olveczky2005, TumerBrainard2007} & измерено \\
12 & Выбор из двух учится при $\tau_e=1\,\text{с}$ и $\delta=R-\bar R$; не учится при $\tau_e=50\ms$; без вычитания веса растут без предела, а при большой скорости обучения сеть может закрепить худший выбор & \texttt{models/dofa\_learning.py}, $\eta=0.05$, $500$ попыток, $6$ зёрен. $\tau_e=1\,\text{с}$: доля выборов $A$ $0.65$--$0.79\to0.96$--$1.00$, веса $0.85$--$1.33$. $\tau_e=50\ms$: $0.38$--$0.55$, веса не меняются. $\delta=R$: вес победителя $860$--$1190$. $\delta=R$, $\eta=0.5$, $3$ зерна: одно закрепилось на $B$ ($0.04\to0$). Скрипт печатает все четыре случая; пересчёт совпал с главой. & \texttt{models/\allowbreak dofa\_\allowbreak learning.py} & модель \\
13 & Время обучения по одному общему сигналу растёт пропорционально числу шумящих нейронов (утверждение~\ref{prop:broadcastcost}) & Кривые обучения линейной сети: возмущение нейронов медленнее точного градиента во столько раз, сколько выходных нейронов, возмущение весов --- ещё во столько раз, сколько входов. & \cite{Werfel2005} & теория \\
14 & Выходы \nt{ДОФА} идут прежде всего к областям выбора действий & Полная реконструкция аксонов отдельных нигростриатных \nt{ДОФА}-нейронов крысы: ветви одной клетки покрывают $0.45$--$5.7\,\%$ объёма стриатума (в среднем $2.7\,\%$), вне стриатума ветвей почти нет. & \cite{Matsuda2009} & измерено \\
15 & Кора учится предсказывать свой вход, ошибка считается на месте & Обзор: в сенсорной коре есть ответы на рассогласование ожидаемого и пришедшего входа. Что это общее правило обучения коры, не доказано. & \cite{Keller2018} & гипотеза \\
16 & Дофамин ведёт себя как ошибка~\eqref{eq:ctd}: всплеск, перенос на предвестник, провал (рис.~\ref{fig:ctdcases}) & \nt{ДОФА}-нейроны обезьяны: всплеск на неожиданный сок; после обучения --- всплеск на предвестник и нет ответа на обещанный сок; провал, когда сок не пришёл. Непрерывная форма~\eqref{eq:ctd} --- теория. & \cite{Schultz1997, Doya2000} & измерено \\
17 & Схема, вычисляющая $\dd V/\dd t$ и вычитающая ожидание & Мыши, запись и оптогенетика в вентральной покрышке: \nt{ДОФА}-нейроны вычитают ожидание из награды; соседние \nt{ГАМК}-нейроны тормозят их, когда награда ожидается, и их возбуждение или торможение меняет ошибку. Как вычисляется производная, не показано. & \cite{Eshel2015arithmetic} & гипотеза \\
18 & \nt{ДОФА}-нейрон --- ритмоводитель; провалы мельче всплесков & В срезе \nt{ДОФА}-нейроны разряжаются спонтанно и очень регулярно, на медленном пейсмекерном потенциале. У обезьяны частота количественно следует положительной ошибке, а отрицательную передаёт лишь слабо. & \cite{GraceOnn1989, Bayer2005} & измерено \\
19 & Актёр и критик (рис.~\ref{fig:actorcritic}) & Алгоритм --- из теории обучения с подкреплением. У людей (фМРТ) брюшной стриатум ведёт себя как критик и при выборе, и без него, спинной --- только когда надо выбирать действия. & \cite{SuttonBarto2018, ODoherty2004actor} & гипотеза \\
20 & Знак $\delta$ в правиле перевёрнут у нейронов со вторым семейством рецепторов & Срезы стриатума: у нейронов с D1-рецепторами спаривание даёт усиление, которому нужен D1; у нейронов с D2 --- ослабление, которому нужен D2. & \cite{Shen2008} & измерено \\
\bottomrule
\end{longtable}}

\section{Глава~\ref{sec:dofaalt}. Другие теории \nt{DOPA}}

Каждая из расширенных картин опирается на свои прямые измерения дофамина (разброс точек переворота, ответы на неприятное и новое, медленный рост, ответы на смену вкуса), но формулы, которые их объясняют, --- теории, и между двумя из них опыты пока противоречат друг другу.

{\footnotesize
\setlength{\tabcolsep}{3pt}\renewcommand{\arraystretch}{1.15}
\begin{longtable}{>{\raggedright}p{0.5cm}>{\raggedright}p{4.0cm}>{\raggedright}p{5.6cm}>{\raggedright}p{2.3cm}>{\raggedright\arraybackslash}p{1.7cm}}
\toprule
№ & Утверждение & Опыт и что измерено & Источник & Статус \\
\midrule\endfirsthead
\toprule
№ & Утверждение & Опыт и что измерено & Источник & Статус \\
\midrule\endhead
1 & Асимметричное правило~\eqref{eq:asymmetric} сходится к точке~\eqref{eq:expectile}; при $\kappa=1/2$ это среднее, для капли с вероятностью $1/2$ --- $V=\kappa$ (рис.~\ref{fig:expectiles}) & Точка~\eqref{eq:expectile} --- экспектиль, решение задачи наименьших квадратов с разными весами положительных и отрицательных отклонений; для примера главы вывод в самой главе. & \cite{NeweyPowell1987}; вывод в главе & теория \\
2 & У разных \nt{DOPA}-нейронов разные точки переворота, и они упорядочены по асимметрии (утверждение~\ref{prop:expectile}) & Мыши, оптогенетически помеченные \nt{DOPA}-нейроны вентральной покрышки, награды разного размера: точка, где всплеск сменяется провалом, у разных нейронов разная; у нейронов с большей асимметрией ответов она выше; по ответам группы восстанавливается форма распределения наград. Что это главная функция разброса --- гипотеза. & \cite{Dabney2020} & измерено \\
3 & Часть \nt{DOPA}-нейронов отвечает на неприятное всплеском & Обезьяны: одни \nt{DOPA}-нейроны возбуждаются на предвестник награды и тормозятся на предвестник струи воздуха в глаз, другие возбуждаются на оба; группы расположены в разных частях среднего мозга. & \cite{Matsumoto2009, BrombergMartin2010} & измерено \\
4 & Модуль ошибки или новизна задаёт скорость обучения & Прямого опыта, где сигнал важности у \nt{DOPA}-нейронов меняет скорость обучения, в цитируемых работах нет; обзор предлагает роль сигнала «внимание, важно». & \cite{BrombergMartin2010} & гипотеза \\
5 & Отдельный провод для оси опасности & Мыши: \nt{DOPA}-нейроны, идущие к заднему концу стриатума, отвечают на новые и сильные раздражители, а не на награду; их разрушение ослабляет избегание нового предмета. & \cite{Menegas2018} & измерено \\
6 & Оптимальное время действия $\tau_a^*=\sqrt{C/\bar R}$~\eqref{eq:vigor} & Минимум суммы $C/\tau_a+\bar R\tau_a$; вывод в главе. В исходной модели скорость действий задаёт цена времени, равная средней награде. & \cite{Niv2007}; вывод в главе & теория \\
7 & Медленный уровень дофамина несёт $\bar R$ и задаёт скорость действий (утверждение~\ref{prop:multiplex}) & Крысы, микродиализ и вольтамперометрия в прилежащем ядре: уровень дофамина от минуты к минуте следует за частотой наград и за скоростью действий. У людей L-DOPA усиливает зависимость времени реакции от средней частоты наград. Что медленный уровень равен именно $\bar R$, не доказано. & \cite{Hamid2016, Beierholm2013vigor} & гипотеза \\
8 & Медленный уровень складывается из двух источников: выброс меняется у выходов без изменения частоты импульсов, но медленный рост бывает и в самой частоте & Крысы, одна задача: выброс в прилежащем ядре следует за меняющимся ожиданием награды, а частота импульсов опознанных \nt{DOPA}-нейронов вентральной покрышки --- нет. Мыши в виртуальном коридоре (строка~10): плавный рост при приближении к награде есть и в импульсах опознанных \nt{DOPA}-нейронов. & \cite{Mohebi2019, Kim2020} & измерено \\
9 & Дофамин плавно растёт, пока животное приближается к далёкой награде & Крысы в лабиринте, вольтамперометрия в стриатуме: концентрация дофамина нарастает за секунды по мере приближения к цели, крутизна зависит от расстояния и размера награды. & \cite{Howe2013ramp, Hamid2016} & измерено \\
10 & Рост --- ошибка $\delta$ при уточнении оценки~\eqref{eq:ramp}; скачок в коридоре даёт всплеск (рис.~\ref{fig:ramp}) & Формула и число $\delta=0.75\,V$ в секунду --- вывод в главе. Опыт: мыши в виртуальном коридоре, мгновенный перенос ближе к цели; импульсы тел, кальций тел и выходов и концентрация в стриатуме отвечают как ошибка, а не как ожидание $V$. Какое из двух объяснений верно на всех выходах, обсуждается. & \cite{Kim2020} & измерено \\
11 & Взгляд назад: дофамин сообщает меру причинной связи~\eqref{eq:retro} & Мыши, выброс дофамина в прилежащем ядре в опытах, где эта теория и ошибка~\eqref{eq:ctd} предсказывают разное: в ряде опытов выброс следовал первой. & \cite{Jeong2022} & гипотеза \\
12 & Ответ дофамина при обучении постепенно сдвигается от награды к предвестнику & Мыши, сигнал выходов \nt{DOPA}-нейронов в брюшном стриатуме по ходу обучения: всплеск постепенно переходит с момента награды на момент предвестника, как требует обучение по~\eqref{eq:ctd}. & \cite{Amo2022} & измерено \\
13 & \nt{DOPA}-нейроны отвечают на смену вкуса награды при той же ценности & Крысы, запись нейронов вентральной покрышки: замена сока другим, равным по ценности, вызывает ответ, хотя ошибка~\eqref{eq:ctd} равна нулю. & \cite{Takahashi2017} & измерено \\
14 & Искусственные всплески учат связи двух нейтральных раздражителей & Крысы, опыт с предварительным сочетанием двух раздражителей без награды: оптогенетическое включение \nt{DOPA}-нейронов на переходе от первого ко второму создаёт связь, выключение мешает её выучить. & \cite{Sharpe2017} & измерено \\
15 & Вектор ошибок по признакам~\eqref{eq:vectortd} & Теория ошибки предсказания признаков; прямой проверки векторной формы нет. & \cite{Gardner2018} & гипотеза \\
16 & Разные \nt{DOPA}-нейроны в одной задаче несут разные величины & Мыши в виртуальном коридоре, двухфотонная съёмка \nt{DOPA}-нейронов: одни связаны с наградой, другие со скоростью и поворотом, третьи с предвестниками и точностью выбора. & \cite{Engelhard2019} & измерено \\
17 & Жгут вместо провода (утверждение~\ref{prop:bundle}) & Сводка строк~2--16: каждое разложение измерено по отдельности; единой картины, что все они --- части одного сигнала, опытом нет. & --- & гипотеза \\
\bottomrule
\end{longtable}}

\section{Глава~\ref{sec:nora}. \nt{НОРА}}

Устройство \nt{НОРА}-системы, её два режима, связь со зрачком и перевёрнутое U в поведении измерены прямо; то, что усиление переключает режим выбора и работает как обратная температура, --- выводы главы; выгода от подстройки усиления --- расчёт по модели книги; «\nt{НОРА} --- ручка поиска» и «\nt{НОРА} --- прерывание» --- гипотезы.

{\footnotesize
\setlength{\tabcolsep}{3pt}\renewcommand{\arraystretch}{1.15}
\begin{longtable}{>{\raggedright}p{0.5cm}>{\raggedright}p{4.0cm}>{\raggedright}p{5.6cm}>{\raggedright}p{2.3cm}>{\raggedright\arraybackslash}p{1.7cm}}
\toprule
№ & Утверждение & Опыт и что измерено & Источник & Статус \\
\midrule\endfirsthead
\toprule
№ & Утверждение & Опыт и что измерено & Источник & Статус \\
\midrule\endhead
1 & \nt{НОРА}-нейронов у человека несколько десятков тысяч, почти все в одной группе & Серийные срезы ствола мозга человека, окраска на тирозингидроксилазу: $53\,900$ клеток в голубом пятне и ещё $6\,260$ в соседнем подъядре. & \cite{Baker1989LC} & измерено \\
2 & Выходы расходятся почти по всему мозгу, но части группы отчасти обслуживают разные области & Вирусно-генетическое мечение у мышей: \nt{НОРА}-нейроны голубого пятна получают входы из широких областей мозга и шлют выходы в широкие области мозга и спинного мозга. У крыс: группа, идущая к миндалине, усиливает выучивание страха, а отдельная группа, идущая к префронтальной коре, --- его угашение. & \cite{Schwarz2015LC, Uematsu2017LC, Poe2020} & измерено \\
3 & Фоновый режим: ровные импульсы с частотой от долей до нескольких в секунду, меняющейся с бодрствованием & Крысы, запись нейронов голубого пятна в цикле сна и бодрствования: частота наибольшая в бодрствовании, ниже в медленном сне, почти ноль в быстром; изменения частоты опережают смену состояния. & \cite{AstonJonesBloom1981} & измерено \\
4 & Отрывистый режим: короткая пачка на важное для задачи событие, примерно в момент решения & Обезьяны, задача найти редкую цель: все нейроны голубого пятна отвечают пачкой только на цель, через $\sim90\ms$ после неё и за $\sim200\ms$ до ответа рукой; при плохой работе пачки слабее. В задаче с выбором из двух пачка точнее привязана к ответу, чем к стимулу, и есть и при верных, и при ошибочных ответах. & \cite{AstonJones1994vigilance, Clayton2004LC} & измерено \\
5 & Диаметр зрачка --- контрольная точка \nt{НОРА} & Обезьяны: импульсы голубого пятна, вплоть до отдельных импульсов одной клетки, предсказывают расширение зрачка; похожая, но более поздняя связь есть в холмиках и поясной коре. Мыши: колебания зрачка следуют за активностью и норадреналиновых, и холинергических выходов в коре --- зрачок отражает не одну \nt{НОРА}. & \cite{Joshi2016, Reimer2016} & измерено \\
6 & \nt{НОРА} меняет крутизну передаточной характеристики~\eqref{eq:gain} & Бодрствующие обезьяны, слуховая кора, ионофорез норадреналина: фоновая активность нейронов подавляется сильнее, чем ответ на голоса сородичей, --- отношение сигнала к шуму растёт. Мультипликативная сигмоида~\eqref{eq:gain} --- упрощение, взятое из сетевой модели. & \cite{Foote1975NE, ServanSchreiber1990} & измерено \\
7 & Усиление повышает $d'$ только против шума после усилителя~\eqref{eq:gainsnr} (утверждение~\ref{prop:gainnoise}) & Вывод в главе; в сетевой модели усиление улучшает обнаружение сигнала. Опыта, который отделял бы шум до и после усилителя и проверял бы это различие, нет. & вывод в главе; \cite{ServanSchreiber1990} & теория \\
8 & Граница $g\beta=1$ переключает сеть выбора из «размышляет» в «решила»~\eqref{eq:wtagain} (утверждение~\ref{prop:gainwta}, рис.~\ref{fig:pitchfork}) & Линейный анализ двух тормозящих друг друга групп; вывод в главе. Прямых измерений этого порога в мозге нет. & вывод в главе & теория \\
9 & Коэффициент усиления --- обратная температура выбора~\eqref{eq:softmax} & При логистическом шуме после усилителя вероятность выбора --- softmax с $T=\sigma/g$; вывод в главе. & вывод в главе & теория \\
10 & \nt{НОРА} задаёт, сколько искать & Люди: перед выбором наугад базовый зрачок шире, чем перед выбором лучшего известного; у людей с более широким зрачком склонность искать выше. Крысы: переход в режим случайного выбора управляется входом \nt{НОРА} в переднюю поясную кору. Направление: больше тонической \nt{НОРА} --- больше поиска, то есть меньше действующее $g$; отождествление \nt{НОРА} с $g$ держится на различии фонового и отрывистого режимов. & \cite{JepmaNieuwenhuis2011, Tervo2014stochastic, AstonJones2005} & гипотеза \\
11 & Награда зависит от $g$ как перевёрнутое U; лучшее $g$ меньше в быстро меняющемся мире (утверждение~\ref{prop:explore}, рис.~\ref{fig:invertedU}) & \texttt{models/nora\_explore.py}, $60$ прогонов по $4000$ попыток. Перемена раз в $1000$ попыток: лучшее $g=16$, доля $0.976$; раз в $20$: $g=8$, доля $0.886$; при $g=0.5$ --- $0.77$. Пересчёт совпал с главой. & \texttt{models/\allowbreak nora\_\allowbreak explore.py} & модель \\
12 & Перевёрнутое U в поведении: лучше всего при умеренной фоновой активности с чёткими пачками & Обезьяны, та же задача, что в строке~4: в периоды хорошей работы фоновая частота умеренная, а пачки на цель чёткие; при высокой фоновой частоте пачек нет и ложных ответов больше. Изменения фоновой и вызванной активности следуют за колебаниями качества работы. & \cite{Usher1999LC, AstonJones2005} & измерено \\
13 & \nt{НОРА} сообщает неожиданную неопределённость, \nt{АЦЕТ} --- ожидаемую & Теория байесовского вывода с двумя видами неопределённости; прямого опыта, разделяющего эти сигналы по медиаторам, в цитируемых работах нет. & \cite{YuDayan2005} & гипотеза \\
14 & После резкой перемены люди учатся быстрее, и зрачок расширяется & Люди, задача предсказания с внезапными переменами: короткие изменения зрачка следуют за оценкой вероятности перемены и вместе с базовым зрачком предсказывают, насколько сильно новые данные меняют оценку (скорость обучения); изменение зрачка, не связанное со светом и задачей, меняет эту скорость. Что зрачок здесь показывает именно \nt{НОРА}, --- вывод по косвенным данным строки~5. & \cite{Nassar2012} & измерено \\
15 & Отрывистая пачка работает как прерывание: сеть сбрасывает состояние & Прямого опыта, где пачка \nt{НОРА} сбрасывает состояние сети, нет; измерены только пачки на важные события (строка~4). & \cite{BouretSara2005} & гипотеза \\
16 & Подстройка $g$ или скорости обучения по неожиданности почти не лучше лучших постоянных настроек & \texttt{models/nora\_explore.py}, мир с эпохами по $1000$ попыток (перемена раз в $1000$ и раз в $20$). Лучшее постоянное $g=8$: $0.925$; подстройка $g$: до $0.927$; подстройка скорости обучения: до $0.926$; постоянная скорость $0.4$: $0.928$. Пересчёт совпал с главой. & \texttt{models/\allowbreak nora\_\allowbreak explore.py} & модель \\
\bottomrule
\end{longtable}}

\section{Глава~\ref{sec:acet}. Добавляем \nt{АЦЕТ}}

Три действия ацетилхолина измерены на срезах и в живом мозге, конфликт записи и чтения показан на модели книги, а то, что \nt{АЦЕТ} служит линией разрешения записи и сообщает ожидаемую неопределённость, --- гипотеза с частичной опорой на опыт.

{\footnotesize
\setlength{\tabcolsep}{3pt}\renewcommand{\arraystretch}{1.15}
\begin{longtable}{>{\raggedright}p{0.5cm}>{\raggedright}p{4.0cm}>{\raggedright}p{5.6cm}>{\raggedright}p{2.3cm}>{\raggedright\arraybackslash}p{1.7cm}}
\toprule
№ & Утверждение & Опыт и что измерено & Источник & Статус \\
\midrule\endfirsthead
\toprule
№ & Утверждение & Опыт и что измерено & Источник & Статус \\
\midrule\endhead
1 & \nt{АЦЕТ}-нейронов мозга немного, они собраны в нескольких группах, а выходы расходятся по всей коре: широковещательный канал & Окраска антителами к ферменту синтеза ацетилхолина и ретроградная метка у крысы: кора, гиппокамп, миндалина, обонятельная луковица и таламус получают ацетилхолин в основном от шести компактных групп клеток базального переднего мозга и ствола & \cite{Mesulam1983} & измерено \\
2 & Уровень ацетилхолина в коре высок при бодрствовании и низок в медленном сне & Микродиализ в коре и гиппокампе свободно движущихся кошек с записью ЭЭГ: выброс ацетилхолина в спокойном бодрствовании на $\sim100\%$, в активном --- на $\sim175\%$ выше, чем в медленном сне; в быстром сне в коре --- как при бодрствовании & \cite{Marrosu1995,Hasselmo1999} & измерено \\
3 & Смысл сигнала задаёт рецептор: у ацетилхолина есть быстрые рецепторы-каналы и медленные, запускающие реакции в клетке (утверждение~\ref{prop:receptor}) & Обзор измерений: действие ацетилхолина на возбудимость, передачу и пластичность зависит от места выброса, подтипа рецептора (никотиновые --- ионные каналы, мускариновые --- через G-белки) и клетки-получателя & \cite{Picciotto2012} & измерено \\
4 & Первое действие: ослабить внутренние связи, почти не трогая входы извне (множитель $1-a$ в~\eqref{eq:acetnet}) & Срезы обонятельной коры крысы: при $100$\,мкМ агонистов ацетилхолина потенциалы внутренних связей (слой Ib) падают более чем на $60\%$, внешнего входа (слой Ia) --- менее чем на $15\%$, в той же клетке; подавление пресинаптическое. Срезы гиппокампа (CA1): $89\%$ против $40\%$ для внутреннего и входного путей & \cite{HasselmoBower1992,HasselmoSchnell1994} & измерено \\
5 & Второе действие: усилить вход (множитель $\frac{1+a}{2}$) & Срезы неокортекса: никотиновые рецепторы усиливают синапсы входа из таламуса и не действуют на внутрикорковые; мускариновые ослабляют оба типа. Возбудимость нейронов ацетилхолин повышает (обзор) & \cite{Gil1997,Hasselmo2006} & измерено \\
6 & Третье действие: облегчить изменение весов (скорость $\eta a$) & Крыса: электрическая стимуляция базального ядра (источника ацетилхолина для коры) в паре с тоном вызывает у взрослого животного сильную перестройку карты слуховой коры под предъявляемый звук & \cite{KilgardMerzenich1998,Hasselmo2006} & измерено \\
7 & Правило Хебба для редких образов во втором уравнении~\eqref{eq:acetnet} даёт ассоциативную память, дописывающую образ по части & Расчёт ёмкости сети с разреженными образами и ковариационным правилом: при малой доле активных нейронов $p$ сеть хранит заметно больше образов, чем при $p=1/2$ & \cite{Tsodyks1988} & теория \\
8 & Запись при низком $a$ портит память: сеть записывает вспомненное, а не увиденное; при $a=0.1$ порча не слабее, чем при $a=0$ & \texttt{models/\allowbreak acet\_memory.py}: $200$ нейронов, пять образов по $20$, новый образ делит с одним из них $10$ нейронов, $10$ прогонов. При $a=0$ новый образ не запомнен, старый тянет $5.9$ чужих нейронов из $10$; при $a=0.1$ новый образ вспоминается ($0.97$), но оба сливаются: новый тянет $7.3$ чужих нейрона, старый --- $7.9$; с $a=0.2$ --- меньше одного & вывод в главе & модель \\
9 & Утверждение~\ref{prop:writeread}: нет уровня $a$, при котором одинаково хороши запись и чтение (рис.~\ref{fig:acetsim}) & Тот же скрипт, скорость обучения $\eta a$: новый образ запоминается за одно предъявление целиком при $a\ge0.9$, на $0.8$ при $a=0.75$, не запоминается при $a\le0.5$. Чтение по половине образа: $1.0$ при $a\le0.2$, $0.96$ при $a=0.3$, $0.04$ при $a=0.4$ & вывод в главе & модель \\
10 & В мозге один широковещательный провод --- \nt{АЦЕТ} --- переключает запись и чтение & Идея взята из моделей, построенных по измерениям на срезах. Самый прямой опыт --- у человека: блокатор мускариновых рецепторов скополамин ухудшает заучивание новых пар слов, сильнее всего пересекающихся со старыми, но не вспоминание уже выученных пар. Прямого измерения двух режимов сети нет & \cite{HasselmoAndersonBower1992,Atri2004} & гипотеза \\
11 & Фильтр~\eqref{eq:kalman} оптимален; вес входа и скорость обучения --- одна величина $P/R$, в установившемся режиме $P=\sqrt{qR}$ и память $\sqrt{R/q}$ (утверждение~\ref{prop:kalman}) & Опыта не требует: оптимальность фильтра Калмана--Бьюси --- классический результат теории оценивания; установившийся режим выведен в главе & \cite{KalmanBucy1961} & теория \\
12 & \nt{АЦЕТ} сообщает сети величину, подобную $P$, --- ожидаемую неопределённость & Предложено в вероятностной модели внимания. Прямого измерения того, что уровень ацетилхолина следует за неопределённостью оценки, нет & \cite{YuDayan2005} & гипотеза \\
13 & Часть \nt{АЦЕТ}-нейронов отвечает на награду и наказание за два десятка миллисекунд, тем сильнее, чем неожиданнее подкрепление & Мышь, оптогенетически опознанные холинергические нейроны базального переднего мозга: ответ на награду и наказание через $18\pm3$\,мс, величина растёт с неожиданностью подкрепления, ответы схожи у нейронов двух ядер & \cite{Hangya2015} & измерено \\
14 & \nt{НОРА} замечает неожиданную перемену мира, \nt{АЦЕТ} держит сеть в режиме записи & Разделение труда предложено в модели. Косвенно: у человека диаметр зрачка, связанный с \nt{НОРА}, отслеживает перемены и скорость обучения. Прямой проверки пары \nt{НОРА}--\nt{АЦЕТ} нет & \cite{YuDayan2005,Nassar2012} & гипотеза \\
15 & В медленном сне сеть в режиме чтения проигрывает записанное днём, и эти повторы переписывают его в долговременную память & Записи ансамблей нейронов гиппокампа крысы: клетки, работавшие вместе днём, чаще срабатывают вместе в последующем медленном сне и в том же порядке --- измерено. Для перезаписи: подавление рябей гиппокампа после обучения ухудшает пространственную память, а удлинение спонтанных рябей её улучшает; что причина --- низкий уровень \nt{АЦЕТ}, не доказано & \cite{WilsonMcNaughton1994,Skaggs1996,Girardeau2009,FernandezRuiz2019,Hasselmo1999} & гипотеза \\
\bottomrule
\end{longtable}}

\section{Глава~\ref{sec:synthesis}. Вся машина}

Глава-сводка новых опытов не приводит: каждая её строка опирается на главу, где утверждение разобрано, и на те же источники; твёрдо установлены шина \nt{ГЛУТ} и управление потоком через \nt{ГАМК}, а роли широковещательных каналов и их совместная работа --- в основном гипотеза.

{\footnotesize
\setlength{\tabcolsep}{3pt}\renewcommand{\arraystretch}{1.15}
\begin{longtable}{>{\raggedright}p{0.5cm}>{\raggedright}p{4.0cm}>{\raggedright}p{5.6cm}>{\raggedright}p{2.3cm}>{\raggedright\arraybackslash}p{1.7cm}}
\toprule
№ & Утверждение & Опыт и что измерено & Источник & Статус \\
\midrule\endfirsthead
\toprule
№ & Утверждение & Опыт и что измерено & Источник & Статус \\
\midrule\endhead
1 & На $8.6\cdot10^{10}$ нейронов приходится всего четыре рода управляющих сигналов~\eqref{eq:machine} & Подсчёт ядер клеток в гомогенате мозга (изотропный фракционатор): у взрослого мужчины $86.1\pm8.1$ млрд нейронов & \cite{Azevedo2009} & измерено \\
2 & Утверждение~\ref{prop:architecture}, первые два слоя: данные идут по адресной шине \nt{ГЛУТ}, знак связи задаёт отправитель, поток направляют и нормируют \nt{ГАМК}-нейроны (главы~\ref{sec:neurontypes}, \ref{sec:gaba}) & Один главный медиатор на нейрон (обзор измерений); прямое торможение сужает окно, в котором пирамидный нейрон складывает входы; деление на общую активность измерено во многих областях & \cite{Strata1999,Pouille2001,Carandini2012} & измерено \\
3 & Элементы ненадёжны: синапс теряет большую часть импульсов (табл.~\ref{tab:computer}, глава~\ref{sec:code}) & Срезы гиппокампа, минимальная стимуляция: больше половины импульсов не вызывают ответа в CA1; сбои порога и проведения по аксону исключены, причина --- вероятностный выброс медиатора & \cite{AllenStevens1994} & измерено \\
4 & Мозг работает на $\sim20$ ваттах, в среднем около импульса в секунду на нейрон (глава~\ref{sec:code}); такой машины инженеры пока не построили & Оценка~\eqref{eq:energy} из измеренных затрат: около $10^9$ молекул АТФ на импульс вместе с работой синапсов (кора грызунов); подробная оценка для коры даёт ещё меньшую частоту. Состояние техники --- по обзору & \cite{Attwell2001,Lennie2003,MehonicKenyon2022} & теория \\
5 & Обучение большинства синапсов --- произведение местного следа и одного общего числа $\delta$ (второе уравнение~\eqref{eq:machine}, глава~\ref{sec:dofa}) & \nt{ДОФА}-нейроны обезьяны отвечают на ошибку предсказания награды; в срезах полосатого тела дофамин увеличивает шипик, только если приходит через $0.3$--$2$\,с после глутаматного входа --- след измерен & \cite{Schultz1997,Yagishita2014} & измерено \\
6 & Отдельный провод ошибки к каждому выходу есть лишь в цепи обучения движениям (глава~\ref{sec:motorlearning}) & Мозжечок: совпадение сигнала лазящего волокна со входом ослабляет этот вход; у обезьяны сложные импульсы клеток Пуркинье несут направление ошибки движения и вызывают поправку & \cite{Ito2001,Herzfeld2018} & измерено \\
7 & Каждый широковещательный канал --- жгут из нескольких линий (утверждение~\ref{prop:bundle}) & Для \nt{ДОФА}: в одной задаче разные нейроны несут награду, движение и признаки стимула; быстрая ошибка и медленный уровень дофамина различны. Для \nt{СЕРО}, \nt{НОРА}, \nt{АЦЕТ} такая раскладка изучена хуже & \cite{Engelhard2019,Hamid2016,Mohebi2019} & измерено \\
8 & \nt{СЕРО}: регулятор уровня и, по гипотезе, горизонт $\tau_\gamma$ (глава~\ref{sec:serocomputer}) & Самоторможение: агонисты собственных рецепторов 5-HT$_{1A}$ тормозят серотониновые нейроны шва --- измерено. Горизонт предложен в теории; сильнейший опыт --- оптогенетическая активация этих нейронов у мыши удлиняет ожидание отложенной награды & \cite{Sprouse1987,Doya2002,Miyazaki2014} & гипотеза \\
9 & \nt{НОРА}: усиление $g$ выбирает между поиском и решением (глава~\ref{sec:nora}) & Рост отношения сигнала к шуму: у бодрствующей обезьяны норадреналин в слуховой коре подавляет фоновую активность сильнее, чем ответ на голоса сородичей --- измерено; мультипликативное усиление --- модель; роль в выборе между поиском и решением --- теория & \cite{Foote1975NE,ServanSchreiber1990,AstonJones2005} & гипотеза \\
10 & \nt{АЦЕТ}: переключает запись и чтение (глава~\ref{sec:acet}) & Три действия (ослабление внутренних связей, усиление входа, облегчение пластичности) измерены; переключение режимов косвенно поддержано опытом со скополамином у человека & \cite{HasselmoBower1992,Gil1997,Atri2004} & гипотеза \\
11 & Формула~\eqref{eq:machine} описывает машину целиком & Это сводка моделей глав~\ref{sec:glutcomputer}--\ref{sec:acet}, каждая часть опирается на свою главу; как целое опытом не проверялась & вывод в главе & гипотеза \\
12 & Ручки работают согласованно без общего управления: ошибка, неожиданность, запись, возврат (рис.~\ref{fig:timeline}) & Опыта нет: схема качественная. Отдельные звенья --- теория разделения труда \nt{НОРА} и \nt{АЦЕТ} и связь зрачка со скоростью обучения у человека & \cite{YuDayan2005,Nassar2012} & гипотеза \\
13 & Обучение по одному общему сигналу медленнее, чем больше нейронов влияют на результат (утверждение~\ref{prop:broadcastcost}) & Расчёт кривых обучения линейной сети: обучение по возмущению выходов медленнее градиентного спуска во столько раз, сколько выходов & \cite{Werfel2005} & теория \\
14 & Как кора настраивает многослойные цепи, неизвестно; кандидаты --- пачки импульсов, вычисления в ветвях нейрона, самообучение предсказанием & Пачки как носитель ошибки --- модель; повторить ответ подробной модели нейрона коры смогла лишь сеть в $5$--$8$ слоёв --- модель; кальциевые импульсы в ветвях нейронов коры человека, дающие исключающее ИЛИ, --- измерено в срезах; самообучение --- обзор свидетельств & \cite{Payeur2021,Beniaguev2021,Gidon2020,Keller2018} & гипотеза \\
\bottomrule
\end{longtable}}

\section{Глава~\ref{sec:sensors}. Входы машины}

Устройство датчиков измерено прямо --- записями тока одиночных клеток, колебаний улитки и подсчётом клеток сетчатки; предел чувствительности и формула дробового шума следуют из физики, а то, что коды датчиков подстроены под наибольшую информацию, --- гипотеза с сильной опорой.

{\footnotesize
\setlength{\tabcolsep}{3pt}\renewcommand{\arraystretch}{1.15}
\begin{longtable}{>{\raggedright}p{0.5cm}>{\raggedright}p{4.0cm}>{\raggedright}p{5.6cm}>{\raggedright}p{2.3cm}>{\raggedright\arraybackslash}p{1.7cm}}
\toprule
№ & Утверждение & Опыт и что измерено & Источник & Статус \\
\midrule\endfirsthead
\toprule
№ & Утверждение & Опыт и что измерено & Источник & Статус \\
\midrule\endhead
1 & Датчик --- преобразователь: рецептор даёт ток, а выход чувствительных клеток глаза и уха --- глутамат на шину \nt{ГЛУТ}; первые клетки импульсов не выдают (рис.~\ref{fig:transducer}) & Палочки и колбочки черепахи выбрасывают возбуждающую аминокислоту: её ловят глутаматные каналы на отсечённом кусочке мембраны. Внутренние волосковые клетки улитки крысы: токи в окончаниях нервных волокон идут через глутаматные AMPA-рецепторы, частота выбросов растёт с плавной деполяризацией клетки до $150$\,Гц & \cite{CopenhagenJahr1989,GlowatzkiFuchs2002} & измерено \\
2 & Датчик света в темноте отвечает на один фотон током около пикоампера & Всасывающий электрод на наружном сегменте палочки жабы: ответы на тусклые вспышки квантованы, событие $\approx1$\,пА ($5\%$ темнового тока) с разбросом $0.2$\,пА, эффективность $0.5$. Человек сообщает о попадании одного фотона чаще случайного & \cite{Baylor1979,Tinsley2016} & измерено \\
3 & Датчик звука замечает смещение меньше нанометра & Метод Мёссбауэра, основная мембрана улитки морской свинки на месте $16$--$19$\,кГц: на пороге ответа слухового нерва скорость мембраны около $0.04$\,мм/с, то есть смещение $\approx0.35$\,нм (пересчёт наш). Мерили мембрану, а не пучок датчика & \cite{Sellick1982,Hudspeth1989} & измерено \\
4 & Точность в темноте ограничена дробовым шумом фотонов~\eqref{eq:photonnoise}; при слабом свете накопление длиннее & Формула следует из пуассоновской статистики. В палочке шум при тусклом свете совпадает с суммой независимых квантовых событий; на ярком фоне ответ на вспышку меньше и короче. Число «десятые доли секунды» здесь отдельно не проверено & \cite{Baylor1979} & теория \\
5 & Диапазон входа --- десять порядков для света и двенадцать для звука, а частота импульсов --- меньше трёх порядков & Для звука: ответ основной мембраны на характеристической частоте растёт с крутизной до $0.2$\,дБ/дБ в полосе $40$--$80$\,дБ, сжатие видно и выше $100$\,дБ. Сами числа порядков --- стандартные оценки, в главе без источника & \cite{Ruggero1997} & измерено \\
6 & Ответ датчика~\eqref{eq:nakarushton}, а $\sigma$ следует за средней яркостью, сдвигая кривую по логарифмической оси (рис.~\ref{fig:adaptation}) & Формула впервые подобрана к S-потенциалам сетчатки рыбы. Колбочки черепахи: ответы подчиняются $V=I V_m/(I+\sigma)$ и охватывают $\sim3.5$ порядка яркости; после $2$--$3$ минут фона кривая сдвигается по горизонтали, сохраняя ширину. Связь с делением на общую активность --- обзор & \cite{NakaRushton1966,NormannPerlman1979,Carandini2012} & измерено \\
7 & Датчики передают перемены, а их коды подстроены под наибольшую информацию на импульс (утверждение~\ref{prop:sensors}) & Принцип предложен теоретически. Сильнейшее свидетельство: у мухи характеристика «контраст $\to$ ответ» нейронов первого слоя совпадает с накопленным распределением контрастов природных сцен --- так достигается наибольшая информация & \cite{Barlow1961,Laughlin1981} & гипотеза \\
8 & Включающие и выключающие датчики --- двухпроводный код; событийная камера повторяет его в кремнии & Обзор опытов: включающий и выключающий каналы сетчатки разделяются уже на первом синапсе и выключаются по отдельности. Камера: $128\times128$ пикселей без кадров, диапазон $120$\,дБ, задержка $15$\,мкс & \cite{Schiller1992,Lichtsteiner2008,Gallego2022} & измерено \\
9 & В глазу $\sim10^8$ датчиков света и $\sim10^6$ выходных нейронов: сжатие в $\sim100$ раз & Подсчёт на целых препаратах сетчатки человека: в среднем $92$ млн палочек и $4.6$ млн колбочек; выходных (ганглиозных) клеток от $0.7$ до $1.5$ млн в разных глазах & \cite{Curcio1990,CurcioAllen1990} & измерено \\
10 & У мыши выходных каналов глаза больше тридцати & Мышь: двухфотонная кальциевая съёмка более $11\,000$ выходных клеток и кластеризация ответов --- заметно больше $30$ функциональных типов. Для человека число не измерено & \cite{Baden2016} & измерено \\
11 & Глаз морской свинки передаёт $\sim875\,000$ бит/с; пересчёт на глаз человека даёт $\sim10^7$ бит/с & Морская свинка, многоэлектродная матрица, движение в природных сценах: $6$--$13$ бит/с на клетку, $\sim875\,000$ бит/с на $\sim10^5$ выходных клеток. Для человека ($\sim10^6$ клеток) $\sim10^7$ бит/с --- пересчёт авторов & \cite{Koch2006} & измерено \\
12 & Ухо --- банк полосовых фильтров с почти логарифмической картой частот (рис.~\ref{fig:filterbank}) & Место наибольшего колебания основной мембраны зависит от частоты; функция «частота--место» почти экспоненциальна, одной формой описывает улитки человека и живых животных & \cite{Greenwood1990,Robles2001} & измерено \\
13 & Резонаторы активны: часть датчиков усиливает слабые колебания в тысячи раз по амплитуде ($66$--$76$\,дБ), с ростом громкости усиление падает & Лазерная виброметрия основания улитки шиншиллы: при слабых тонах на характеристической частоте усиление $66$--$76$\,дБ относительно стремечка; смерть снижает чувствительность на $60$--$81$\,дБ (в $10^3$--$10^4$ раз по амплитуде) и убирает сжатие. Источник силы --- наружные волосковые клетки & \cite{Ruggero1997,Robles2001} & измерено \\
14 & Усилитель уха стоит у границы самовозбуждения & Расчёт: у точки бифуркации Хопфа неизбежны сжатие диапазона, острая настройка и комбинационные тоны --- все известные нелинейности слуха. Что улитка настроена именно так, прямо не доказано & \cite{Eguiluz2000} & гипотеза \\
15 & На низких частотах импульсы идут в фазе со звуком (у морской свинки привязка слабеет от $\sim600$\,Гц и заметна до $\sim3.5$\,кГц), и по ним находится направление; на высоких остаётся код местом & Слуховой нерв морской свинки: привязка к фазе начинает падать около $600$\,Гц и исчезает выше $3.5$\,кГц; у кошки и обезьяны граница примерно на октаву выше. Разница времени прихода в два уха --- обзор & \cite{PalmerRussell1986,Grothe2010} & измерено \\
16 & Остальные датчики (табл.~\ref{tab:sensors}): у обоняния $\sim400$ типов рецепторов, по одному на нейрон, и комбинаторный код; вкус частично через АТФ и серотонин; осязание --- быстро и медленно привыкающие датчики; тепло и холод --- отдельные & Мышь: один рецептор узнаёт много запахов, один запах --- много рецепторов, разные запахи --- разные наборы. Без рецепторов АТФ нерв вкуса не отвечает; вкусовые почки выбрасывают серотонин. Запись одиночных волокон кожи руки человека: четыре типа по скорости привыкания. Холодовой канал отличен от тепловых & \cite{Mombaerts2004,Malnic1999,Finger2005,Huang2005,JohanssonVallbo1979,McKemy2002} & измерено \\
\bottomrule
\end{longtable}}

\section{Глава~\ref{sec:recognition}. Распознавание}

Скорость распознавания, устройство первых ступеней, линейное чтение ответа и обучение на непрерывности времени измерены на человеке и обезьяне; сведение всей цепочки к двум операциям и обучение на видео проверены на моделях книги, а объяснения иллюзий --- от хорошо обоснованных до спорных.

{\footnotesize
\setlength{\tabcolsep}{3pt}\renewcommand{\arraystretch}{1.15}
\begin{longtable}{>{\raggedright}p{0.5cm}>{\raggedright}p{4.0cm}>{\raggedright}p{5.6cm}>{\raggedright}p{2.3cm}>{\raggedright\arraybackslash}p{1.7cm}}
\toprule
№ & Утверждение & Опыт и что измерено & Источник & Статус \\
\midrule\endfirsthead
\toprule
№ & Утверждение & Опыт и что измерено & Источник & Статус \\
\midrule\endhead
1 & Сравнение с образцом по пикселям при сдвиге угадывает не лучше монетки; пара ступеней~\eqref{eq:scstage} узнаёт фигуру в любом месте & \texttt{models/\allowbreak recognition\_models.py}, «сетчатка» из $24$ точек, $4000$ проб: образец --- $0.49$, $0.51$, $0.52$ при доле перевёрнутых точек $0$, $0.1$, $0.2$; пара $S$ и $C$ --- $1.00$, $0.86$, $0.70$ & вывод в главе & модель \\
2 & Животное на картинке узнаётся за $\sim150$\,мс, один проход по десятку ступеней по $10$--$20$\,мс & Человек, задача «есть животное или нет» на фотографиях, показанных на $20$\,мс: вызванные потенциалы для двух ответов расходятся через $\sim150$\,мс. Обезьяна: время первого ответа растёт от ступени к ступени (V1 раньше всех, затем V2 и V4). Число ступеней и «ноль или один импульс на ступень» --- вывод из этих чисел & \cite{Thorpe1996,Schmolesky1998} & измерено \\
3 & Ступень $S$ --- И по частям, ступень $C$ --- наибольшее по положениям (утверждение~\ref{prop:conveyor}) & Кошка, первая зрительная кора: простые клетки отвечают на край определённой ориентации в определённом месте, сложные --- на ту же ориентацию в большей области. Внутриклеточная запись сложных клеток: ответ на две полоски у многих клеток близок к большему из двух ответов, в среднем ближе всего к операции max. Наибольшее через взаимное торможение --- утверждение~\ref{prop:wta}, деление на общую активность --- обзор & \cite{HubelWiesel1962,Lampl2004,Carandini2012} & измерено \\
4 & Выше по цепочке сочетания крупнее и сложнее, ответ всё меньше зависит от положения & Обезьяна, упрощение стимулов до «критического признака»: клетки V4 и задней нижневисочной коры отвечают на сложные сочетания формы, цвета и текстуры при малом поле; в передней нижневисочной коре поле больше. Чередование $S$ и $C$ в модели воспроизводит эту картину & \cite{KobatakeTanaka1994,Riesenhuber1999} & измерено \\
5 & Свёрточные сети повторяют это чередование и лучше всех предсказывают ответы верхних ступеней; предмет читается линейно по частотам группы нейронов & Сеть, обученная распознавать, без подгонки к мозгу: верхний слой предсказывает ответы нейронов нижневисочной коры обезьяны, средние --- ответы V4. Линейный классификатор по активности $\sim100$ случайно выбранных клеток за окна от $12.5$\,мс узнаёт предмет и категорию при разном положении и размере & \cite{Fukushima1980,Yamins2014,Hung2005} & измерено \\
6 & Правило Хебба со следом~\eqref{eq:traceinvariance} учит инвариантности на видео без подписей, если след не короче $\sim20$\,мс & Идея медленных признаков и правила со следом --- теория. \texttt{models/\allowbreak recognition\_models.py}, два $C$-нейрона, $16$ входов, $10$ прогонов: пауза между предметами $0.3$\,с. При $\tau_{\text{сл}}=5$\,мс совпадение с фигурой $0.52$ в среднем, при $10$\,мс --- $0.56$; при $20$, $50$ и $200$\,мс --- $1.00$ во всех прогонах (при $30$\,мс один прогон из десяти ошибается) & \cite{Foldiak1991,WiskottSejnowski2002} & модель \\
7 & В мозге инвариантность учится на непрерывности времени & Обезьяна: предмет незаметно подменяли во время скачка глаза в определённое место поля; позиционная инвариантность нейронов нижневисочной коры менялась в соответствии с подменой, значимо уже через час. Что это главное правило обучения зрения --- гипотеза & \cite{LiDiCarlo2008} & измерено \\
8 & Движение: детекторы направления, затем та же пара И/ИЛИ по большим областям & Детекторы направления в сетчатке кролика; в области MT обезьяны все $110$ записанных нейронов избирательны к направлению, ответ на движение сильнее, чем в V1 & \cite{Barlow1965,Albright1984} & измерено \\
9 & Верхние ступени посылают вниз предсказание, вверх идёт ошибка & Модель объясняет внеклассические эффекты полей первой зрительной коры; отдельные наблюдения согласуются (обзор), как общее устройство конвейера не доказано & \cite{RaoBallard1999,Keller2018} & гипотеза \\
10 & Трудные картинки требуют обратных связей и нескольких кругов & Обезьяна: для «трудных» картинок решение в нижневисочной коре складывается на $\sim30$\,мс позже; эти поздние ответы плохо предсказывает прямая сеть и лучше --- сети с обратными связями или более глубокие & \cite{Kar2019} & измерено \\
11 & Незаметная подделка пикселей обманывает сеть; зрение человека гораздо устойчивее, но не неуязвимо & На сетях: малое специально подобранное возмущение меняет ответ; пример «панда $\to$ гиббон» --- из второй работы. У человека при коротком показе подделки, хорошо переносимые между сетями, сдвигают выбор & \cite{Szegedy2014,Goodfellow2015,Elsayed2018} & измерено \\
12 & Иллюзии ожидания: ненадёжный вход уступает ожиданию --- бледное движется «медленнее», платье видят по-разному (раздел~\ref{sec:illusions}) & Решётки меньшего контраста кажутся движущимися медленнее. Опрос $1401$ человека: платье видят в трёх устойчивых вариантах, подсказка об освещении переключает цвет; у $\sim13\,000$ наблюдателей решает предположение об освещении. Байесовская модель с ожиданием медленных движений объясняет ряд иллюзий движения & \cite{Thompson1982,LaferSousa2015,Wallisch2017,Weiss2002,Purves2011} & гипотеза \\
13 & Регулировка усиления: водопад, наклон и контраст; решётка Германа --- не торможение соседей & Детекторы направления сетчатки кролика при долгом движении снижают частоту, а после остановки падают ниже фона. Иллюзия наклона воспроизводится моделью с делением на активность соседей. Изменения решётки, не меняющие ответ выходных нейронов сетчатки, убирают пятна --- объяснение соседним торможением отвергнуто & \cite{BarlowHill1963,Schwartz2009,SchillerCarvey2005} & измерено \\
14 & Компенсация задержек: движущийся предмет кажется впереди вспышки & Иллюзия измерена. Психофизика противоречит и продлению движения вперёд, и разнице задержек; предложено объяснение задним числом --- по событиям $\sim80$\,мс после вспышки. Спор не решён & \cite{Nijhawan1994,EaglemanSejnowski2000} & гипотеза \\
15 & Неподвижные «змеи» вращаются: разница задержек~\eqref{eq:latency} читается детекторами направления & Иллюзия работает и без цвета; пары соседних элементов порождают её поодиночке; нейроны коры обезьяны, избирательные к направлению, дают направленный ответ на эти неподвижные пары. У человека вращение запускают микроскачки глаз и моргания & \cite{Conway2005,OteroMillan2012} & измерено \\
16 & Двойственные картинки: взаимное торможение, усталость и шум; смены вызывает главным образом шум & Модель конкурирующих групп нейронов: смены в ней вызывает \emph{шум}, без шума их нет; она объясняет рост частоты смен с силой стимула. Усталость здесь играет меньшую роль & \cite{MorenoBote2007} & гипотеза \\
17 & Часть иллюзий --- следствие задачи, которой учится машина & Сеть, обученная предсказывать следующий кадр видео, «видит» вращение неподвижных колец и не видит его на контрольных картинках; сети для очистки и повышения резкости изображений поддаются иллюзиям цвета и яркости, как человек & \cite{Watanabe2018,GomezVilla2020} & измерено \\
\bottomrule
\end{longtable}}

\section{Глава~\ref{sec:muscles}. Выход на мышцы}

Устройство выхода --- двигательные единицы, запас надёжности синапса, включение по размеру, пара мышц, петля растяжения и генераторы ритма --- измерено прямо; условие устойчивости петли с задержкой --- теорема; внутренняя модель тела --- хорошо обоснованная гипотеза; числа рисунков --- расчёт по модели книги.

{\footnotesize
\setlength{\tabcolsep}{3pt}\renewcommand{\arraystretch}{1.15}
\begin{longtable}{>{\raggedright}p{0.5cm}>{\raggedright}p{4.0cm}>{\raggedright}p{5.6cm}>{\raggedright}p{2.3cm}>{\raggedright\arraybackslash}p{1.7cm}}
\toprule
№ & Утверждение & Опыт и что измерено & Источник & Статус \\
\midrule\endfirsthead
\toprule
№ & Утверждение & Опыт и что измерено & Источник & Статус \\
\midrule\endhead
1 & Двигательная единица: один \nt{АЦЕТ}-нейрон управляет группой волокон, от нескольких сотен до пары тысяч; в мышцах тонкой подстройки единицы мельче & Мышцы человека, подсчёт двигательных аксонов и мышечных волокон: в среднем около $340$ волокон на нейрон в межкостной мышце кисти, около $410$ в плечелучевой мышце, около $1930$ в медиальной икроножной. Точного числа для самых мелких единиц глава не приводит. & \cite{Feinstein1955mu} & измерено \\
2 & Синапс на мышце выбрасывает в несколько раз больше медиатора, чем нужно для срабатывания волокна & Обзор измерений на нервно-мышечных синапсах: запас надёжности (отношение выброшенного медиатора к пороговому) у взрослых млекопитающих обычно $3$--$5$; у человека запас держится больше на складках мембраны получателя, чем на величине выброса. & \cite{WoodSlater2001} & измерено \\
3 & Подёргивание~\eqref{eq:twitch} нарастает за время $T_c$ порядка $50\ms$ & Одиночные двигательные единицы кисти человека при произвольном усилии, сила усреднена по импульсам единицы: время нарастания $30$--$100\ms$, у $80\,\%$ единиц меньше $70\ms$. Функция $(t/T_c)e^{1-t/T_c}$ --- аппроксимация из модели пула единиц. & \cite{MilnerBrown1973rec, Fuglevand1993} & измерено \\
4 & Сумма подёргиваний~\eqref{eq:forcesum}: при $5$, $15$ и $40$ имп/с сила $0.2$--$1.1F_1$, $1.8$--$2.2F_1$ и около $5.4F_1$ (рис.~\ref{fig:tetanus}) & Расчёт по \texttt{models/\allowbreak muscle\_\allowbreak models.py}, $T_c=50\ms$: $0.21$--$1.10$, $1.76$--$2.19$ и $5.28$--$5.49\,F_1$ на последних $0.3\,\text{с}$. Линейное сложение --- упрощение: насыщение настоящей мышцы в модели нет. & \texttt{models/\allowbreak muscle\_\allowbreak models.py} & модель \\
5 & Единицы включаются по размеру; самые слабые в сотню раз слабее самых сильных & Вентральные корешки кошки: при нарастающем рефлекторном входе первыми разряжаются нейроны с мелкими импульсами в корешке, то есть мелкие. Межкостная мышца кисти человека: сила подёргивания единиц от $0.1$ до $10$~г и растёт почти линейно с силой, при которой единица включается. & \cite{Henneman1957, MilnerBrown1973rec} & измерено \\
6 & Шаг силы пропорционален силе, $\Delta F/F\approx q-1$~\eqref{eq:sizeprinciple}: плавающая точка & Выведено в главе из геометрической прогрессии сил. Согласуется с опытом: при больших усилиях на то же приращение силы включается гораздо меньше новых единиц. & вывод в главе; \cite{MilnerBrown1973rec} & теория \\
7 & Разброс силы растёт пропорционально самой силе & Произвольное изометрическое усилие мышцы большого пальца: стандартное отклонение силы линейно растёт со средней силой; при том же усилии, вызванном электрической стимуляцией, разброс постоянен. Модель пула: линейный рост даёт включение единиц по порядку силы. & \cite{Jones2002sdn} & измерено \\
8 & Плавность движений --- следствие шума, зависящего от сигнала & Модель: траектория, минимизирующая разброс конечной точки при таком шуме, воспроизводит измеренные траектории глаза и руки. Прямого опыта, отличающего эту причину от других, нет. & \cite{HarrisWolpert1998} & гипотеза \\
9 & Разность сил пары мышц задаёт момент, сумма --- жёсткость~\eqref{eq:antagonists} & Теория управления импедансом совместным напряжением. Рука человека: жёсткость каждого сустава, измеренная малыми толчками, примерно пропорциональна моменту в нём; разные соотношения совместного напряжения меняют форму и направление эллипса жёсткости кисти. & \cite{Hogan1984, GomiOsu1998stiff} & измерено \\
10 & Датчик растяжения возбуждает свой \nt{АЦЕТ}-нейрон и через тормозящего посредника выключает противника (рис.~\ref{fig:antagonists}) & Спинной мозг кошки: одиночный посредник, возбуждаемый волокнами датчиков растяжения, даёт моносинаптические тормозные потенциалы в \nt{АЦЕТ}-нейронах мышцы-противника, синаптическая задержка $0.28$--$0.42\ms$. Медиатор посредника (глицин) в этом опыте не определялся. & \cite{Jankowska1972Ia} & измерено \\
11 & Задержка петли: спинномозговой $25$--$30\ms$, через кору в полтора--три раза больше, через глаз $100$--$200\ms$ & Рука человека, толчок по суставу: короткий ответ мышц через $20$--$45\ms$, длинный $45$--$105\ms$, произвольный $120$--$180\ms$. Смещение видимого положения пальца во время движения: поправка через $160\ms$ в среднем. & \cite{Pruszynski2008rapid, SaundersKnill2003vis} & измерено \\
12 & $\dd x/\dd t=-Gx(t-d)$ устойчиво при $Gd<\pi/2$~\eqref{eq:delaystable}; на границе частота $1/(4d)$ & Теорема о корнях уравнения с запаздыванием. Проверка по \texttt{models/\allowbreak muscle\_\allowbreak models.py}, $d=30\ms$: к $3\,\text{с}$ амплитуда $3\cdot10^{-6}$ при $G=40$, $0.13$ при $G=50$, $38$ при $G=55$ в секунду (граница $52$). & \cite{Hayes1950} & теория \\
13 & Физиологическая дрожь $8$--$12$~Гц; петля растяжения --- один из её источников & Обзор измерений: мелкая дрожь в диапазоне около $10$~Гц есть у здоровых людей; её происхождение смешанное --- центральные колебания, свойства разрядов единиц, механический резонанс и резонанс рефлекторной петли, и в разных условиях преобладает разное. & \cite{McAuleyMarsden2000} & гипотеза \\
14 & Мозг предсказывает результат команды по внутренней модели тела & Человек держит груз щипком и двигает рукой: при инерционной, вязкой и упругой нагрузке сила сжатия меняется одновременно с нагрузкой, а не с запаздыванием петли, то есть нагрузка предсказана. Обзор сводит такие опыты; прямого наблюдения самой модели в нейронах нет. & \cite{FlanaganWing1997grip, WolpertGhahramani2000} & гипотеза \\
15 & Ритм ходьбы, дыхания, жевания порождают местные сети при постоянном входе & Обзор опытов: изолированная нервная система (спинной мозг, ганглии) выдаёт ритмический двигательный узор без ритмического входа и без обратной связи от мышц. & \cite{MarderBucher2001} & измерено \\
16 & Взаимное торможение с усталостью~\eqref{eq:halfcentre} даёт ритм с периодом $0.84\,\text{с}\approx2\tau_a$ (рис.~\ref{fig:halfcentre}) & Теорема: сеть со взаимным торможением и адаптацией без устойчивого состояния при постоянном входе обязательно колеблется. Расчёт по \texttt{models/\allowbreak muscle\_\allowbreak models.py} ($I=1$, $\beta=2$, $b=2$, $\tau=20\ms$, $\tau_a=0.4\,\text{с}$): период $0.84\,\text{с}$. Что настоящие генераторы устроены именно так, не показано. & \cite{Matsuoka1985osc} & модель \\
\bottomrule
\end{longtable}}

\section{Глава~\ref{sec:pain}. Боль}

Датчики боли, два провода, ворота в спинном мозге, нисходящая ручка громкости, сенсибилизация и захват внимания измерены прямо; формулы ворот и сенсибилизации --- упрощённые модели книги; правило, по которому машина выбирает громкость тревоги, --- гипотеза.

{\footnotesize
\setlength{\tabcolsep}{3pt}\renewcommand{\arraystretch}{1.15}
\begin{longtable}{>{\raggedright}p{0.5cm}>{\raggedright}p{4.0cm}>{\raggedright}p{5.6cm}>{\raggedright}p{2.3cm}>{\raggedright\arraybackslash}p{1.7cm}}
\toprule
№ & Утверждение & Опыт и что измерено & Источник & Статус \\
\midrule\endfirsthead
\toprule
№ & Утверждение & Опыт и что измерено & Источник & Статус \\
\midrule\endhead
1 & Высокий порог: датчики боли срабатывают при нагреве выше примерно $43^\circ$ & Запись одиночных волокон. Кожа обезьяны: медианный порог медленных (C) датчиков $41^\circ$, быстрых датчиков второго типа $46^\circ$, первого типа --- выше $53^\circ$. Кожа человека: C-датчики, отвечающие и на давление, $40.7^\circ$, не отвечающие на давление --- $48^\circ$. Число $43^\circ$ --- середина этого диапазона, а не единый порог. & \cite{Treede1995heat, Weidner1999cfib} & измерено \\
2 & После повреждения датчики не глохнут, а становятся чувствительнее & Лёгкий ожог кожи ($50^\circ$, $100\,\text{с}$): через $5$--$10$~мин порог боли у людей и порог C-датчиков у обезьяны снижены на $2$--$6^\circ$; у других кожных датчиков такого сдвига нет. Слабое привыкание прямо не сравнивалось; у быстрых датчиков второго типа ответ после сильного нагрева, наоборот, подавлен. & \cite{LaMotte1982hyper, Treede1995heat} & измерено \\
3 & Два провода: быстрый $5$--$30$~м/с, медленный $0.5$--$2$~м/с; время прихода $t=L/v$~\eqref{eq:paintime} & Скорость проведения: быстрые датчики боли обезьяны в среднем $15$ и $25$~м/с (два типа); C-датчики человека $0.8$--$1.0$~м/с. При $L=1$~м это десятки миллисекунд и около секунды. & \cite{Treede1995heat, Weidner1999cfib} & измерено \\
4 & Боль приходит дважды: сначала быстрая и точная, потом тупая и долгая & Время реакции на нагрев предплечья человека: от $1100$ до $400\ms$ с ростом температуры; сильный нагрев волосистой кожи даёт двойную боль, ладони --- нет. Первую боль могут нести только волокна быстрее $6$~м/с --- по латентности ей соответствуют быстрые датчики второго типа, которых в ладони нет. Раздел ролей («где» и «насколько плохо») --- толкование. & \cite{CampbellLaMotte1983first, Treede1995heat} & измерено \\
5 & Ворота: выходной нейрон спинного мозга получает боль и касание, тормозящий посредник касанием возбуждается (рис.~\ref{fig:gate}) & Схема ворот предложена как теория. Мыши, генетическая метка и удаление групп нейронов: возбуждающие нейроны с соматостатином нужны для механической боли; тормозящие нейроны с динорфином не дают входу от датчиков касания дойти до них --- без этих нейронов лёгкое касание вызывает боль. & \cite{MelzackWall1965, Duan2014} & измерено \\
6 & Касание приглушает боль & Человек, лазерные импульсы боли на руке: одновременное касание снижает оценку боли и способность различать энергию импульсов; эффект слабеет с расстоянием между точкой касания и точкой боли внутри одного кожного сегмента. & \cite{Mancini2014touch} & измерено \\
7 & Сильная боль сама выключает тормозящего посредника, и ворота распахиваются & Это часть исходной схемы ворот. В приведённых опытах прямо не показано: работа на мышах описывает, как ворота не пускают касание в канал боли, а не выключение посредника болью. & \cite{MelzackWall1965} & гипотеза \\
8 & Ворота~\eqref{eq:gate}: порог $0.18$ без касания, $0.86$ с касанием, $0.11$ при $g=2$; при $\nu=1$ касание снижает выход с $1$ до $0.25$, при $\nu=2$ не действует; касание само не проходит (рис.~\ref{fig:gatecurves}) & Расчёт по \texttt{models/\allowbreak pain\_\allowbreak gate.py} при весах из текста воспроизводит все эти числа. & \texttt{models/\allowbreak pain\_\allowbreak gate.py} & модель \\
9 & Нисходящие пути из ствола меняют усиление ворот в обе стороны & Крыса, продолговатый мозг, запись при нагреве хвоста: одни нейроны разряжаются перед отдёргиванием («вкл»), другие замолкают («выкл»); слабая стимуляция этой области подавляет рефлекс. Обзор: те же цепи и облегчают, и тормозят передачу боли, опиоиды действуют через них. & \cite{Fields1983rvm, Fields2004} & измерено \\
10 & Ожидание облегчения приглушает боль через опиоидный канал & Люди с острой болью: у тех, кому ожидание облегчения снизило боль, блокада опиоидных рецепторов вернула боль к уровню остальных; у остальных блокада боль не меняла. & \cite{Levine1978} & измерено \\
11 & Громкость тревоги мозг выбирает по выгоде прерывания & Опыта, в котором измерено правило выбора громкости, нет. & --- & гипотеза \\
12 & Сенсибилизация: после повреждения выход ворот растёт, порог падает, отчасти за счёт спинного мозга & Крыса: после повреждения кожи порог сгибательного рефлекса падает, а ответ растёт; электрофизиологический анализ показывает, что часть роста возникает в самом спинном мозге. Время спада «от часов до дней» в этой работе не измерялось. & \cite{Woolf1983} & измерено \\
13 & Сенсибилизация~\eqref{eq:sensitization} --- положительная обратная связь; при сильной петле тревога может защёлкнуться после заживления & Следует из модели~\eqref{eq:sensitization} (задача в конце главы). Опыта, показывающего, что стойкая боль после заживления --- защёлкнувшаяся петля именно такого вида, нет. & вывод в главе & гипотеза \\
14 & Боль --- отрицательная награда, и обучение на ней следует ошибке временных разностей & Томограф, люди, цепочки предвестников боли: активность вентрального стриатума и передней островковой коры совпадает с сигналами обучения по временным разностям. Обезьяна: часть \nt{ДОФА}-нейронов возбуждается от неприятного обдува и его предвестника, другая часть тормозится. & \cite{Seymour2004, Matsumoto2009} & измерено \\
15 & Боль прерывает текущую работу & Люди, электрический болевой стимул и звуковые пробы: ответ на пробу через $100\ms$ после начала боли заметно замедлен, через $1.5\,\text{с}$ разницы с контролем уже нет. Обзор сводит такие опыты в модель захвата внимания. & \cite{Crombez1996disrupt, EcclestonCrombez1999} & измерено \\
16 & Отдёргивание замыкается в спинном мозге & Крыса: нейроны глубоких слоёв заднего рога повторяют пространственный узор рефлекса отдёргивания отдельных мышц; их не удаётся возбудить антидромно из шейного отдела, то есть это спинальные посредники. & \cite{Schouenborg1995wr} & измерено \\
\bottomrule
\end{longtable}}

\section{Глава~\ref{sec:motorlearning}. Обучение движениям}

Правило наименьших квадратов и выгода отдельного провода ошибки --- теоремы; устройство схемы с проводом ошибки, ослабление при совпадении, подстройка следа под задержку, последействие и спонтанный возврат навыка измерены; то, что вся схема работает как обучение с учителем и строит внутреннюю модель, --- ведущая гипотеза; числа двухпроцессной модели --- расчёт по модели книги.

{\footnotesize
\setlength{\tabcolsep}{3pt}\renewcommand{\arraystretch}{1.15}
\begin{longtable}{>{\raggedright}p{0.5cm}>{\raggedright}p{4.0cm}>{\raggedright}p{5.6cm}>{\raggedright}p{2.3cm}>{\raggedright\arraybackslash}p{1.7cm}}
\toprule
№ & Утверждение & Опыт и что измерено & Источник & Статус \\
\midrule\endfirsthead
\toprule
№ & Утверждение & Опыт и что измерено & Источник & Статус \\
\midrule\endhead
1 & Правило~\eqref{eq:lms} в среднем идёт по градиенту квадрата ошибки & Результат теории адаптивных фильтров: поправка, пропорциональная входу и ошибке выхода, есть стохастический спуск по градиенту среднего квадрата ошибки. & \cite{WidrowHoff1960} & теория \\
2 & С проводом ошибки на каждый выход скорость не зависит от числа выходов; с одним общим числом --- падает с числом шумящих нейронов (утверждение~\ref{prop:teachers}) & Кривые обучения линейной сети: обучение по случайным возмущениям нейронов медленнее точного градиента во столько раз, сколько возмущаемых нейронов. & \cite{Werfel2005} & теория \\
3 & Схема с проводом ошибки работает как обучение с учителем (утверждение~\ref{prop:errorwire}) & Теории, выведенные из строения схемы. Опыты строк 7--10 согласуются с ними; что вся схема работает именно так, не доказано. & \cite{Marr1969, Albus1971} & гипотеза \\
4 & Расширитель: около $7\cdot10^{10}$ мелких \nt{GLUT}-нейронов, больше, чем во всём остальном мозге & Подсчёт ядер во взвеси растворённой ткани: в мозжечке человека $69\cdot10^9$ нейронов из $86\cdot10^9$ во всём мозге. Стереология: $101\cdot10^9$ мелких нейронов в мозжечке пожилых мужчин. & \cite{Azevedo2009, Andersen1992cereb} & измерено \\
5 & Подъём размерности входа делает нужную зависимость линейной & Теорема о числе разбиений: взвешенная сумма $n$ входов с порогом реализует произвольную разметку примерно $2n$ векторов общего положения. Что мозжечок пользуется именно этим, --- часть гипотезы строки~3. & \cite{Cover1965} & теория \\
6 & Около $3\cdot10^7$ выходных \nt{GABA}-нейронов, у каждого порядка $10^5$ входов & Стереология мозжечка человека: $30.5\cdot10^6$ выходных нейронов. Электронная микроскопия, крыса: около $1.75\cdot10^5$ входов от расширителя на один выходной нейрон. Тормозной медиатор выходных нейронов --- в обзоре. & \cite{Andersen1992cereb, NapperHarvey1988, Ito2001} & измерено \\
7 & Ровно один провод ошибки на выходной нейрон, около раза в секунду, каждый раз мощная вспышка & Обзор записей: каждый выходной нейрон получает один лазящий \nt{GLUT}-вход, который вызывает сложный импульс с частотой порядка $1$~Гц. & \cite{Ito2001} & измерено \\
8 & Вероятность вспышки зависит от направления ошибки движения & Обезьяна, глазодвигательный отдел мозжечка, адаптация саккад: направление зрительной ошибки задаёт вероятность сложного импульса, величина --- его время; вспышка меняет простые импульсы в следующей попытке и сдвигает движение вдоль предпочтительной ошибки клетки. & \cite{Herzfeld2018} & измерено \\
9 & Совпадение вспышки ошибки с активным входом ослабляет вход ($-\eta e_i x_j$) & Обзор опытов: совместная стимуляция входа от расширителя и провода ошибки вызывает долговременное ослабление этого входа; стимуляция одного входа --- нет. & \cite{Ito2001} & измерено \\
10 & Длина следа подстроена под задержку ошибки: при задержке $100\ms$ сильнее всего ослабляется вход, бывший активным за $100\ms$ & Срезы и живые мыши: в отделе, отвечающем за глазодвигательное обучение, ослабление узко настроено на интервал около $120\ms$ между входом и ошибкой --- столько идёт сигнал ошибки в эту задачу; в другом отделе разные клетки настроены на разные интервалы. & \cite{Suvrathan2016} & измерено \\
11 & Схема --- адаптивный фильтр~\eqref{eq:adaptivefilter}, выучивающий внутреннюю модель тела & Модель, выведенная из строения схемы. Прямого опыта, где выученный фильтр измерен как передаточная функция тела, нет. & \cite{Fujita1982} & гипотеза \\
12 & Предсказание вычитает следствия собственных движений, поэтому себя не пощекотать & Люди: собственное касание кажется менее щекотным, чем то же чужое; в томографе соматосенсорная кора отвечает на собственное касание слабее, а мозжечок меньше активен, когда движение касается кожи. Объяснение через вычитание предсказания --- гипотеза. & \cite{Blakemore1998} & гипотеза \\
13 & При внешней силе промах уменьшается, а после её снятия рука промахивается в другую сторону & Люди тянутся рукоятью робота в силовом поле: траектории возвращаются к прямым; после снятия поля последействие --- почти зеркальное отражение первых искажений; оно переносится и в не тренированные части пространства. & \cite{ShadmehrMussaIvaldi1994} & измерено \\
14 & Учатся два процесса~\eqref{eq:tworate}; после обратного возмущения навык возвращается сам & Люди, силовое поле, затем короткое обратное поле и попытки с зажатой ошибкой: поправка сама возвращается к старой; модель из двух процессов это воспроизводит, модель из одного --- нет. & \cite{Smith2006} & измерено \\
15 & Числа рис.~\ref{fig:tworate}: $0.84$ (медленный $0.63$) за $200$ движений; $6$ обратных; быстрый $-0.53$, медленный $0.46$; возврат до $0.37$ за $40$ движений & Расчёт по \texttt{models/\allowbreak motor\_\allowbreak learning.py}, $A_f=0.92$, $B_f=0.1$, $A_s=0.996$, $B_s=0.02$: $0.840$ ($0.634$ и $0.205$), $6$ движений, $-0.531$ и $0.457$, пик $0.371$ на $40$-м движении. & \texttt{models/\allowbreak motor\_\allowbreak learning.py} & модель \\
16 & Два процесса нужны, потому что тело меняется с разными скоростями & Теория: оптимальная оценка при помехах с несколькими временами жизни даёт несколько фильтров с разными постоянными времени и объясняет спонтанный возврат и ускоренное повторное обучение. & \cite{Kording2007} & гипотеза \\
\bottomrule
\end{longtable}}

\section{Глава~\ref{sec:chemoutput}. Химический выход}

Два провода к сердцу и их разная скорость, отрицательная обратная связь в гормональных каскадах, код частотой порций, суточный генератор и его подстройка светом измерены прямо; граница $K<8$ --- теорема теории регулирования, выведенная в главе; пульсации, порождаемые самой обратной связью каскада, --- гипотеза с сильным опытом в её пользу.

{\footnotesize
\setlength{\tabcolsep}{3pt}\renewcommand{\arraystretch}{1.15}
\begin{longtable}{>{\raggedright}p{0.5cm}>{\raggedright}p{4.0cm}>{\raggedright}p{5.6cm}>{\raggedright}p{2.3cm}>{\raggedright\arraybackslash}p{1.7cm}}
\toprule
№ & Утверждение & Опыт и что измерено & Источник & Статус \\
\midrule\endfirsthead
\toprule
№ & Утверждение & Опыт и что измерено & Источник & Статус \\
\midrule\endhead
1 & Ритмоводитель сердца сам бьётся около $100$ раз в минуту; в покое тормоз нажат, и частота обычно порядка $60$--$70$ & Люди, полное выключение обоих проводов к сердцу: собственная частота выше частоты покоя и убывает с возрастом, у взрослых порядка $100$ ударов в минуту. Значит, в покое преобладает тормоз. Частота покоя $60$--$70$ --- типичное значение, отдельно не цитировано. & \cite{JoseCollison1970ihr} & измерено \\
2 & Газ \nt{NORA} и тормоз \nt{ACET} --- фильтры нижних частот, тормоз быстрее~\eqref{eq:gasbrake} & Собака, частотно-модулированная стимуляция блуждающего или симпатического нерва сердца и передаточная функция к частоте предсердий: оба пути --- фильтры нижних частот, у симпатического частота среза много ниже и добавлена чистая задержка около $1.7\,\text{с}$. Характеристики зависят от среднего тонуса, так что линейная формула~\eqref{eq:gasbrake} --- приближение. & \cite{Berger1989} & измерено \\
3 & Тормоз быстр, потому что калиевый канал \nt{ACET} открывает без второго посредника, а \nt{NORA} --- через цепочку реакций & Мембранные лоскуты клеток предсердий: калиевый канал, открываемый ацетилхолином, открывают $\beta\gamma$-субъединицы G-белка, связанного с рецептором, без второго посредника внутри клетки. Путь \nt{NORA} через цАМФ здесь не цитирован. & \cite{Logothetis1987gbg} & измерено \\
4 & Кровь обходит тело за время порядка минуты; \nt{ADRE} из кровотока доходит до всех органов с рецептором & Общая оценка порядка величины; в главе так и сформулирована, источник не цитирован. & --- & оценка \\
5 & Гормональный каскад охвачен отрицательной обратной связью: конечный гормон тормозит первые ступени & Обзор опытов: гормоны коры надпочечников подавляют выброс АКТГ гипофизом и рилизинг-гормона гипоталамусом в трёх временных диапазонах --- быстром, промежуточном и медленном. & \cite{KellerWood1984fb} & измерено \\
6 & Три одинаковые ступени устойчивы при $K<8$~\eqref{eq:cascadestable}, на границе период $2\pi\tau/\sqrt3\approx3.6\tau$ & Выведено в главе: на частоте $\sqrt3/\tau$ три ступени дают сдвиг фазы $180^\circ$ и ослабление в $8$ раз. & вывод в главе & теория \\
7 & Ошибка уровня $1/(1+K)$, поэтому точнее примерно $10\,\%$ такой регулятор не держит (утверждение~\ref{prop:cascade}) & Следствие строки~6 для пропорциональной обратной связи. Настоящая ось имеет обратную связь на нескольких ступенях и в нескольких временных диапазонах (строка~5), так что граница относится к модели~\eqref{eq:cascade}, а не измерена. & вывод в главе & теория \\
8 & При $K=4$ и $7$ уровень устанавливается, при $K=12$ --- ровные пульсации с периодом $3.7$--$3.9\,\tau$ (рис.~\ref{fig:cascade}) & Расчёт по \texttt{models/\allowbreak chem\_\allowbreak models.py}: при $K=12$ последовательные периоды $3.9$, $3.7$, $3.8$, $3.7\,\tau$; при $K=4$ размах в конце счёта $0.003$. & \texttt{models/\allowbreak chem\_\allowbreak models.py} & модель \\
9 & Гормон стресса выходит пульсами около раза в час; пульсы порождает сама обратная связь, без отдельного генератора & Крысы, постоянное вливание рилизинг-гормона: АКТГ и кортикостерон колеблются с физиологической, почти часовой частотой --- ритмичный вход от гипоталамуса для пульсов не нужен. Модель гипофиз--надпочечник с запаздывающей обратной связью даёт такие колебания. & \cite{Walker2012glc, Walker2010} & гипотеза \\
10 & Получатель читает частоту порций, а не уровень & Обезьяны без собственного выброса рилизинг-гормона гонадотропинов: постоянное вливание не восстанавливает выброс гормонов гипофиза, порции раз в час восстанавливают; переход на постоянное вливание снова глушит ответ. Утомление рецептора как фильтр верхних частот --- толкование. & \cite{Belchetz1978} & измерено \\
11 & Разные частоты порций по-разному действуют на разные клетки одной железы & Те же обезьяны: учащение порций до $2$--$5$ в час снижает оба гормона гипофиза; урежение до одной за $3$~часа снижает один (ЛГ) и повышает другой (ФСГ), так что их отношение задаётся частотой. & \cite{Wildt1981gnrh} & измерено \\
12 & Суточный ритм порождает внутриклеточная отрицательная обратная связь с задержкой в несколько часов & Обзор: белки часовых генов с задержкой подавляют транскрипцию собственных генов; петля из транскрипции и трансляции даёт период около суток. & \cite{TakahashiJS2017} & измерено \\
13 & Главный генератор --- группа из десятков тысяч нейронов, синхронизирующая остальное тело & Обзор: надперекрёстное ядро --- главный суточный генератор; отдельные его нейроны в культуре --- самостоятельные генераторы, сеть их синхронизирует; у мыши около $10^4$ нейронов с каждой стороны. & \cite{Welsh2010scn} & измерено \\
14 & Собственный период у человека $24.18$~часа & Люди в тусклом свете при расписании, оторванном от суток: период ритмов мелатонина, температуры и кортизола в среднем $24.18$~часа у молодых и пожилых, с узким разбросом. & \cite{Czeisler1999} & измерено \\
15 & Свет подводит генератор как фазовая автоподстройка~\eqref{eq:pll}; нужен сдвиг около $11$~мин в сутки & Люди, одиночный яркий свет $6.7$~часа в разное время суток: кривая фазового ответа первого типа с размахом $5$~часов --- задержка до минимума температуры тела, опережение после. Уравнение~\eqref{eq:pll} --- стандартная модель; $11$~мин $=0.18$~часа --- арифметика. & \cite{Khalsa2003prc} & измерено \\
16 & Без света подстройки нет, и ритм уходит на собственный период & Полностью слепые люди без восприятия света, живущие в обычном обществе: примерно у половины ритмы мелатонина и кортизола идут со своим периодом, не совпадающим с сутками. & \cite{Sack1992blind} & измерено \\
\bottomrule
\end{longtable}}

\section{Глава~\ref{sec:debugging}. Отладка машины}

Что слышит электрод, маломерность активности, работа интерфейсов на жёстких решётках, совместное обучение мозга и декодера и зрительные точки от стимуляции измерены в рецензируемых опытах; оценки потоков данных --- арифметика главы; о вживлённом людям устройстве Neuralink, насколько удалось проверить, публиковала данные в основном сама компания.

{\footnotesize
\setlength{\tabcolsep}{3pt}\renewcommand{\arraystretch}{1.15}
\begin{longtable}{>{\raggedright}p{0.5cm}>{\raggedright}p{4.0cm}>{\raggedright}p{5.6cm}>{\raggedright}p{2.3cm}>{\raggedright\arraybackslash}p{1.7cm}}
\toprule
№ & Утверждение & Опыт и что измерено & Источник & Статус \\
\midrule\endfirsthead
\toprule
№ & Утверждение & Опыт и что измерено & Источник & Статус \\
\midrule\endhead
1 & Электрод слышит импульсы нейронов в радиусе порядка сотни микрометров, обычно от одного до нескольких; их различают по форме & Крыса, поле CA1, одновременная внутриклеточная и внеклеточная запись: форма внеклеточного импульса повторяет ширину и амплитуду внутриклеточного. Оценка: тетрод мог бы различать $60$--$100$ нейронов, а на деле выделяют около шести. & \cite{Henze2000extra} & измерено \\
2 & В мозге $8.6\cdot10^{10}$ нейронов; щуп слышит примерно одну стомиллионную & Подсчёт ядер во взвеси растворённой ткани: $86\cdot10^9$ нейронов в мозге человека. Доля --- арифметика главы. & \cite{Azevedo2009} & измерено \\
3 & Активность сотен нейронов двигательной коры укладывается в десяток-другой общих переменных & Обзор записей в двигательной коре обезьян: активность популяции описывается немногими скрытыми переменными, общими для многих нейронов. & \cite{Gallego2017} & измерено \\
4 & Шум соседей согласован, и сотня нейронов несёт почти столько же, сколько тысяча (утверждение~\ref{prop:pooling}) & Зрительная кора обезьяны: коэффициент корреляции шума соседних нейронов около $0.1$, и выигрыш от усреднения насыщается при сотне нейронов. Теория ограничивающих информацию корреляций. & \cite{Zohary1994, MorenoBote2014} & измерено \\
5 & Сырой поток $2\cdot10^8$~бит/с против $8\cdot10^4$~бит/с в импульсах~\eqref{eq:probedata} & Арифметика главы по ёмкости потока импульсов~\eqref{eq:spikeentropy}. Параметры оцифровки реальны: кристалл Neuralink --- $19.3$~кГц, $10$~бит на канал. & вывод в главе; \cite{Rieke1997, Musk2019} & теория \\
6 & Импульсы выделяют на кристалле у электродов и передают наружу только события; корпус закрыт кожей, связь и питание без проводов & Система, описанная в статье компании, передаёт полный сырой поток по кабелю USB-C, импульсы выделяет программа в реальном времени вне кристалла; сжатие на борту, питание и передача без проводов названы там планом для клинических устройств. Для вживлённого людям устройства --- только сообщения компании. & \cite{Musk2019}; отчёты Neuralink, рецензируемой статьи нет & измерено (отчёт компании) \\
7 & До $3072$ электродов на $96$ нитях по $32$; нити вставляет робот, обходя сосуды & Статья компании: $3072$ электрода на $96$ нитях по $32$; робот вставляет $6$ нитей ($192$ электрода) в минуту; до $70\,\%$ электродов у крыс с долговременно вживлённой решёткой слышат импульсы. & \cite{Musk2019} & измерено (статья компании) \\
8 & У людей около тысячи каналов на $64$ нитях; часть нитей сместилась; курсор порядка $10$~бит/с & Только сообщения компании. Поиск в PubMed не нашёл рецензируемой статьи с результатами на людях; это согласуется с оговоркой в тексте главы. & отчёты Neuralink; рецензируемой статьи нет & измерено (отчёт компании) \\
9 & Интерфейс на решётке из сотни электродов управляет курсором & Человек с параличом, $96$ электродов в первичной двигательной коре: намерение двигать рукой меняет импульсы через три года после травмы; декодер даёт «нейронный курсор» (почта, телевизор), управление протезом кисти и роботом-рукой. & \cite{Hochberg2006} & измерено \\
10 & Письмо «мысленной рукой»: около $90$ знаков в минуту, меньше $8$~бит/с & Человек с параличом кисти, попытки писать от руки, рекуррентный декодер: $90$ знаков в минуту при точности $94.1\,\%$ в реальном времени, выше $99\,\%$ после автокоррекции. Оценка в битах --- арифметика главы. & \cite{Willett2021} & измерено \\
11 & Попытки говорить переводятся в текст со скоростью около $60$ слов в минуту & Человек, потерявший внятную речь, решётки в коре: $62$ слова в минуту; $9.1\,\%$ ошибок в словах при словаре из $50$ слов, $23.8\,\%$ при словаре из $125\,000$. & \cite{Willett2023} & измерено \\
12 & Интерфейс вынимает малую долю ёмкости: у одного нейрона десятки бит в секунду, у сотни --- тысячи (утверждение~\ref{prop:probe}) & Верхняя граница~\eqref{eq:spikeentropy} --- теория; сравнение со строками~8, 10 --- вывод главы. Что узкое место --- маломерность и незнание кода, а не провод, прямым опытом не проверено. & \cite{Rieke1997} & теория \\
13 & Мозг и декодер учатся друг у друга & Обезьяны управляют курсором; декодер подстраивается в замкнутой петле, а нейроны одновременно переучиваются: точность растёт, навык сохраняется и меньше страдает от движений собственной руки. Совет развести скорости обучения --- инженерное правило, отдельного опыта в главе нет. & \cite{Orsborn2014coad} & измерено \\
14 & Ток через электрод возбуждает окружающие нейроны и провода без разбора типа & Кора мыши и кошки, двухфотонная запись кальция: даже слабый ток возбуждает редкие нейроны, разбросанные на расстояние до миллиметров, через прямое возбуждение аксонов в объёме диаметром в десятки микрометров. То есть возбуждение неизбирательно, но не сплошное. & \cite{Histed2009stim} & измерено \\
15 & Стимуляция зрительной коры зажигает точки; до $16$ точек сразу позволяют узнать некоторые буквы и границы предметов & Слепой человек, $96$ электродов в зрительной коре на $6$ месяцев: порог одного электрода $66.8\pm36.5$~мкА; одновременная стимуляция групп (до $16$ электродов --- предел стимулятора) снижает порог и даёт различимые узоры, некоторые буквы и границы предметов узнаются. & \cite{Fernandez2021} & измерено \\
16 & Второе направление Neuralink --- устройство, подающее сигналы в зрительную кору & Только сообщения компании о разработке; рецензируемых данных о его работе не найдено. & отчёты Neuralink & гипотеза \\
17 & Концентрации широковещательных медиаторов можно мерить химическим датчиком в ткани & Срезы мозга крысы, быстрая циклическая вольтамперометрия на углеродном волокне: измерены выброс и обратный захват серотонина; вещество выходит за пределы синапса. & \cite{Bunin1998} & измерено \\
\bottomrule
\end{longtable}}

\section{Глава~\ref{sec:eetypes}. Нейроны как приборы}

Режимы работы отдельных клеток измерены прямо, током через электрод и записью напряжения; математика интегратора и деления на классы --- классическая теория, а то, что мозг пользуется перестройкой режимов для вычислений, остаётся гипотезой.

{\footnotesize
\setlength{\tabcolsep}{3pt}\renewcommand{\arraystretch}{1.15}
\begin{longtable}{>{\raggedright}p{0.5cm}>{\raggedright}p{4.0cm}>{\raggedright}p{5.6cm}>{\raggedright}p{2.3cm}>{\raggedright\arraybackslash}p{1.7cm}}
\toprule
№ & Утверждение & Опыт и что измерено & Источник & Статус \\
\midrule\endfirsthead
\toprule
№ & Утверждение & Опыт и что измерено & Источник & Статус \\
\midrule\endhead
1 & Резонатор: медленный ток, противодействующий изменению напряжения, делает нейрон полосовым фильтром с максимумом на единицах--десятках герц (раздел~\ref{sec:resonator}) & В срезах в \nt{GLUT}-нейроны вводили синусоидальный ток с плавно растущей частотой и мерили импеданс $|Z(f)|$: максимум на $2$--$5$~Гц при $33^\circ$C (около $7$~Гц при $38^\circ$C); резонанс создают медленные калиевый и катионный токи и усиливает постоянный натриевый ток. Обзор даёт тот же приём для многих классов клеток & \cite{HuVervaekeStorm2002,HutcheonYarom2000} & измерено \\
2 & Медленный ток ведёт себя как индуктивность, вместе с ёмкостью мембраны --- колебательный контур & Подпороговый отклик аксона на ток измерен и описан линеаризованными уравнениями проводимостей; эквивалентная схема содержит «феноменологическую» индуктивность, величина которой зависит от напряжения & \cite{Mauro1970} & теория \\
3 & Постоянная времени группы с обратной связью $\tau_{\text{эфф}}=\tau/(1-w)$~\eqref{eq:perfectint}; при $\tau=0.1$~с и $w=0.995$ --- $20$~с & Выводится в главе из линейного уравнения группы; как механизм удержания положения глаза предложено в теоретической работе & вывод в главе; \cite{Seung1996} & теория \\
4 & Интегратор положения глаза держит вход сетью, а не свойством отдельной клетки & Внутриклеточная запись у рыбы в узле, держащем положение глаза: после быстрого поворота частота нейрона ступенькой меняется и держится; короткий ток в одну клетку вызывает лишь кратковременный сдвиг, а ступеньки напряжения сохраняются и ниже порога --- их поддерживают синаптические входы & \cite{Aksay2001} & измерено \\
5 & Интегратору без утечки нужна постоянная подстройка обучением; при расстройке он либо забывает, либо уходит в насыщение & Рыбу обучали зрительной обратной связью, пропорциональной $\pm$ положению глаза: удержание становилось неустойчивым или с утечкой, постоянная времени нейронов падала больше чем на порядок, до $<1$~с; обычная зрительная обратная связь постепенно возвращала устойчивость & \cite{Major2004} & измерено \\
6 & Бистабильный нейрон: короткий вход включает длительное плато, торможение выключает; свойство включается серотонином (раздел~\ref{sec:bistable}) & Внутриклеточная запись \nt{ACET}-нейронов, управляющих мышцами, у кошки: короткое синаптическое возбуждение переводит нейрон в длительную работу, короткое торможение сбрасывает; у спинальной кошки состояние появляется после введения предшественника серотонина & \cite{Hounsgaard1988} & измерено \\
7 & Два класса: интеграторы (частота растёт от нуля) и резонаторы (при пороге сразу конечная частота) & Классы выведены из теории бифуркаций. Опыт: в срезах коры \nt{GLUT}-нейроны начинают работать со сколь угодно малой частотой (тип~1), быстрые \nt{GABA}-нейроны --- скачком с конечной частоты (тип~2) & \cite{Izhikevich2007,Tateno2004} & измерено \\
8 & Один нейрон --- интегратор или детектор совпадений: торможение укорачивает окно сложения & В срезах гиппокампа окно, в котором возбуждающие входы складываются до порога, короче $2$~мс; его укорачивает прямое торможение, приходящее сразу за возбуждением; в дендритах, где торможение слабее, окно шире & \cite{Pouille2001} & измерено \\
9 & Линия задержки из аксонов (таблица~\ref{tab:eetypes}) & У сипухи измерены задержки в аксонах, идущих к детекторам совпадений: разная длина аксона даёт разную задержку, и детекторы настроены на разность времён прихода звука & \cite{Carr1990} & измерено \\
10 & Ритмоводитель в бодрствовании и молчаливая клетка во сне & \nt{SERO}-нейроны работают почти регулярно днём, замедляются в медленном сне и почти молчат в быстром (подробнее --- глава~\ref{sec:sleep}) & \cite{JacobsAzmitia1992,McGintyHarper1976} & измерено \\
11 & Прибор задают ионные каналы и широковещательные каналы, которые их подстраивают (утверждение~\ref{prop:eetypes}) & Показано на малых сетях: вещества-регуляторы меняют собственные свойства нейронов и силу синапсов, так что одна и та же схема выдаёт разные ритмы & \cite{Marder2012} & измерено \\
12 & Мозг использует перестройку режимов для вычислений (утверждение~\ref{prop:eetypes}) & Прямого опыта, где перестройка режима клетки была бы нужна для решения задачи, нет; есть лишь опыты, где регуляторы меняют выход схемы & \cite{Marder2012} & гипотеза \\
\bottomrule
\end{longtable}}

\section{Глава~\ref{sec:dendrites}. Число дендритов}

Строение классов и числа входов измерены анатомически и электрическими записями; оценка~\eqref{eq:dendriteload} --- простая физика конденсатора, а вычислительное объяснение числа входов опирается на теорию и модели.

{\footnotesize
\setlength{\tabcolsep}{3pt}\renewcommand{\arraystretch}{1.15}
\begin{longtable}{>{\raggedright}p{0.5cm}>{\raggedright}p{4.0cm}>{\raggedright}p{5.6cm}>{\raggedright}p{2.3cm}>{\raggedright\arraybackslash}p{1.7cm}}
\toprule
№ & Утверждение & Опыт и что измерено & Источник & Статус \\
\midrule\endfirsthead
\toprule
№ & Утверждение & Опыт и что измерено & Источник & Статус \\
\midrule\endhead
1 & Удельная ёмкость мембраны $c_m\approx1$~мкФ/см$^2$ & С тела нейрона снимали мембрану на пипетку, подавали ступеньку напряжения и по спаду ёмкостного тока находили полную ёмкость, затем делили на площадь: $0.9$~мкФ/см$^2$ у трёх классов нейронов & \cite{Gentet2000} & измерено \\
2 & Время заряда $t\approx c_mA\Delta V/I$~\eqref{eq:dendriteload}: большому нейрону нужны десятки одновременных входов, маленькому хватает одного крупного синапса & Оценка по закону заряда конденсатора; числа ($20$ и $300$~пФ, $0.1$ и несколько нА) --- порядки величин & вывод в главе & теория \\
3 & Один огромный синапс даёт ток в несколько наноампер и сразу доводит маленький нейрон до порога & Одновременная запись окончания и получателя в срезе: ток одного синапса $2$--$13$~нА, каждый импульс входа вызывает надпороговый ответ, задержка около $0.4$~мс при $36^\circ$C; получатель выдаёт \nt{ГЛИЦ} & \cite{Borst1995,BorstSoria2012} & измерено \\
4 & Ноль дендритов: у чувствительных нейронов осязания и боли импульс идёт от датчика в спинной мозг мимо тела клетки & Строение (аксон, делящийся надвое у тела) установлено анатомически. То, что проведение не требует возбудимости тела, показано на проверенной модели такого нейрона: без возбудимого тела импульс проходит развилку так же надёжно & \cite{AmirDevor2003} & теория \\
5 & Один дендрит: обонятельный нейрон несёт один тип рецептора и передаёт одну входную линию & Обзор опытов с генами рецепторов: каждый обонятельный нейрон выражает один ген рецептора из большого семейства & \cite{Mombaerts2004} & измерено \\
6 & Два дендрита: у детекторов направления на звук входы от левого и правого уха садятся на разные дендриты & Анатомия и записи у млекопитающих, собранные в обзоре: входы от двух ушей разнесены по двум противоположным дендритам & \cite{Grothe2010} & измерено \\
7 & Разнесение входов по двум дендритам делает «И» строже & Показано на модели: насыщение~\eqref{eq:shunt} на одном дендрите не даёт входам с одной стороны дойти до порога, и детекция совпадений улучшается. Прямого опыта, где бы меняли раскладку входов, нет & \cite{AgmonSnir1998} & теория \\
8 & Около $7\cdot10^{10}$ маленьких \nt{ГЛУТ}-нейронов цепи обучения движениям, по $\sim4$ входа у каждого & Подсчёт клеток изотропным фракционированием: в этой части мозга человека около $69$ млрд нейронов. Запись одной такой клетки у живой крысы: на касание приходят пачки дискретных токов от немногих входов, и клетка срабатывает, когда они складываются & \cite{Azevedo2009,Chadderton2004} & измерено \\
9 & Пороговый элемент с $n$ входами разделяет не больше $P_{\max}\approx2n$ случайных образов~\eqref{eq:cover} & Классический результат комбинаторной геометрии. Оговорка: при одних лишь положительных весах и с запасом на шум ёмкость меньше; для выходного нейрона цепи обучения движениям оценено $\sim4\cdot10^4$ соответствий, а не $2\cdot10^5$ & \cite{Cover1965,Brunel2004} & теория \\
10 & Смесители с малым числом входов нужны, чтобы расширить вход; «четыре» близко к оптимуму & Теория: размерность представления, которую даёт слой смесителей, максимальна примерно при том числе входов, которое наблюдается анатомически; исходная гипотеза расширения --- в классических работах & \cite{Marr1969,Albus1971,LitwinKumar2017} & гипотеза \\
11 & Большие деревья: $10^4$--$10^5$ входов & Подсчёт синапсов по электронным снимкам: $\approx1.7\cdot10^5$ синапсов на одном выходном \nt{ГАМК}-нейроне цепи обучения движениям у крысы; около $30\,000$ возбуждающих и $1\,700$ тормозящих входов на \nt{ГЛУТ}-нейроне гиппокампа & \cite{NapperHarvey1988,Megias2001} & измерено \\
12 & Ветви большого дерева --- отдельные подсхемы со своими порогами; нейрон --- маленькая многослойная сеть & Точечная стимуляция двух мест на тонких дендритах: входы на одной ветви складываются по сигмоиде, на разных ветвях --- линейно. В дендритах нейронов коры человека найдены градуальные кальциевые импульсы, позволяющие одной клетке решать линейно неразделимую задачу. Модель: чтобы воспроизвести один нейрон, нужна глубокая сеть & \cite{Polsky2004,Gidon2020,Beniaguev2021} & измерено \\
13 & Дендриты-выходы: каждая ветвь \nt{ГАМК}-нейрона сетчатки --- отдельный детектор направления & Двухфотонная запись кальция в отдельных ветвях: сигнал в ветви избирателен к движению от тела клетки к кончику, а напряжение тела --- нет & \cite{Euler2002} & измерено \\
14 & Число входов следует за задачей (утверждение~\ref{prop:dendrites}) & Строение классов измерено (строки 3--13); что число входов выбрано ради скорости или классификации, показано лишь теорией и моделями для отдельных классов & \cite{LitwinKumar2017,AgmonSnir1998} & гипотеза \\
\bottomrule
\end{longtable}}

\section{Глава~\ref{sec:sleep}. Сон}

Режимы сна, повторы и уменьшение синапсов измерены записями нейронов, микродиализом и электронной микроскопией; расписание сна и медленные качели описаны моделями книги, а назначение сна --- в основном гипотезы.

{\footnotesize
\setlength{\tabcolsep}{3pt}\renewcommand{\arraystretch}{1.15}
\begin{longtable}{>{\raggedright}p{0.5cm}>{\raggedright}p{4.0cm}>{\raggedright}p{5.6cm}>{\raggedright}p{2.3cm}>{\raggedright\arraybackslash}p{1.7cm}}
\toprule
№ & Утверждение & Опыт и что измерено & Источник & Статус \\
\midrule\endfirsthead
\toprule
№ & Утверждение & Опыт и что измерено & Источник & Статус \\
\midrule\endhead
1 & Во сне нейроны коры не молчат; меняется режим работы & Долгие записи нейронов коры у крыс: в медленном сне частота остаётся ненулевой, периоды работы чередуются с периодами тишины; после долгого бодрствования частота выше во всех состояниях, после сна --- ниже & \cite{Vyazovskiy2009} & измерено \\
2 & \nt{ACET} высок днём и в быстром сне, низок в медленном (таблица~\ref{tab:sleepmodes}) & Микродиализ в коре у свободно двигающихся кошек: в бодрствовании выброс ацетилхолина на $100$--$175\%$ выше, чем в медленном сне, а в быстром сне --- примерно на уровне спокойного бодрствования & \cite{Marrosu1995} & измерено \\
3 & \nt{NORA} и \nt{SERO} затихают в медленном сне и почти молчат в быстром & Запись отдельных \nt{NORA}-нейронов у крыс и \nt{SERO}-нейронов у кошек по стадиям сна: частота максимальна в бодрствовании, ниже в медленном сне и почти нулевая в быстром & \cite{AstonJonesBloom1981,McGintyHarper1976} & измерено \\
4 & В быстром сне мышцы выключены торможением через \nt{GABA} и \nt{GLYC} & У крыс мерили активность \nt{ACET}-нейронов жевательной мышцы и подавали блокаторы прямо к ним: паралич снимается, только если заблокированы и \nt{GABA}-рецепторы обоих типов, и \nt{GLYC}-рецепторы & \cite{BrooksPeever2012} & измерено \\
5 & Давление сна $S$ --- конденсатор с двумя порогами~\eqref{eq:sleeppressure}, $\tau_r=18.2$~ч, $\tau_d=4.2$~ч & Модель двух процессов; постоянные времени получены из того, как мощность медленных волн ЭЭГ растёт с бодрствованием и падает во сне, а форма порогов --- из времени спонтанного пробуждения & \cite{Borbely1982,Daan1984} & теория \\
6 & Машина сама выходит на $16.1$~ч бодрствования и $8.0$~ч сна (рис.~\ref{fig:twoprocess}) & Расчёт по~\eqref{eq:sleeppressure} с качающимися порогами, скрипт \texttt{models/sleep\_models.py}: $16.1$~ч бодрствования и $7.9$--$8.0$~ч сна & --- & модель \\
7 & После $40$~ч без сна $S=0.90$, но сон длится лишь $8.8$~ч: долг возвращается глубиной, а не длительностью & Модель: \texttt{models/sleep\_models.py} даёт $S=0.90$ и $8.8$~ч. Опыт: после $40.5$~ч без сна у восьми человек в первую восстановительную ночь мощность медленных волн ЭЭГ значимо выше обычной & \cite{Borbely1981} & модель \\
8 & Физически $S$ --- аденозин & Микродиализ у кошек: внеклеточный аденозин в одной из областей, управляющих бодрствованием, растёт за время бодрствования, растёт дальше при лишении сна и падает во время сна. Что $S$ --- только аденозин, не показано & \cite{PorkkaHeiskanen1997,Dunwiddie2001} & гипотеза \\
9 & В медленном сне кора качается: работа несколько сотен мс --- тишина, реже раза в секунду ($<1$~Гц, под наркозом чаще $0.3$--$0.4$~Гц) & Внутриклеточная запись у кошек под наркозом: деполяризация $0.8$--$1.5$~с и долгая гиперполяризация, ритм $<1$~Гц, чаще всего $0.3$--$0.4$~Гц; это же чередование работы и тишины видно в естественном сне (строка~1) & \cite{Steriade1993,Vyazovskiy2009} & измерено \\
10 & Качели --- группа с обратной связью и усталостью~\eqref{eq:updown}; при $b=5$ период $1.1$~с (значение модели) и работа около $35\%$ времени, при $b=1$ работа ровная (рис.~\ref{fig:updown}) & Расчёт, скрипт \texttt{models/sleep\_models.py}: период $1.11$~с, доля времени работы $0.35$ (среднее по целому числу периодов); при $b=1$ доля работы $1.0$. Период модели короче измеренного (строка~9) & --- & модель \\
11 & \nt{ACET} выключает качели и переводит кору в ровную работу & Стимуляция ядра \nt{ACET}-нейронов у кошки блокировала медленный ритм в $79\%$ нейронов коры и переводила их в ровную работу, в основном за счёт исчезновения долгих гиперполяризаций. Ацетилхолин подавляет медленные калиевые токи, создающие усталость & \cite{SteriadeAmzica1993,McCormick1992} & измерено \\
12 & Вспышки ритма $7$--$15$~Гц порождает петля \nt{GLUT}--\nt{GABA} между корой и релейным узлом & В срезах релейного узла зрения у хорька вспышки создаются ответными пачками релейных клеток на торможение от соседнего \nt{GABA}-ядра, которое они сами возбуждают & \cite{vonKrosigk1993,SteriadeMcCormickSejnowski1993} & измерено \\
13 & Повторы дня идут с перемоткой: в гиппокампе примерно в $20$ раз быстрее, в коре --- в $6$--$7$ раз & В медленном сне повторяющиеся последовательности нейронов гиппокампа идут сжатыми во времени. После бега по дорожке последовательности нейронов места повторялись в том же порядке короткими вспышками по $\sim100$~мс, сжатыми примерно в $20$ раз; в коре сжатие в $6$--$7$ раз & \cite{Nadasdy1999,LeeWilson2002,Euston2007} & измерено \\
14 & Усиление согласованности повторов с корой улучшает память (причинная связь) & У крыс после обучения стимуляцией, привязанной к повторам, усиливали их связь с медленными волнами и вспышками ритма в коре: на следующий день память была высокой, у контрольных --- на уровне случая & \cite{Maingret2016} & измерено \\
15 & Назначение повторов --- перемешанное обучение медленной памяти & Теория двух систем обучения и расчёты на искусственных сетях: последовательное обучение портит старое, перемешанное --- нет. Прямого опыта на мозге, что повторы нужны именно для этого, нет & \cite{McClelland1995} & гипотеза \\
16 & После сна синапсы уменьшаются пропорционально размеру, крупные почти не меняются & Трёхмерная электронная микроскопия $6920$ синапсов коры мыши: площадь контакта после сна на $\sim18\%$ меньше, чем после бодрствования, уменьшение пропорционально размеру, крупные устойчивые синапсы не меняются & \cite{deVivo2017} & измерено \\
17 & Главная задача сна --- перенормировка весов & Гипотеза синаптического гомеостаза; строки 1 и 16 с ней согласуются, но не доказывают, что это главная функция & \cite{TononiCirelli2014} & гипотеза \\
18 & Быстрый сон --- испытание модели мира на стенде & Режим (таблица~\ref{tab:sleepmodes}) измерен; что сеть прогоняет модель мира, --- теоретическое предположение, опыта нет & \cite{Hobson2009} & гипотеза \\
19 & Промывка мозга во сне: вопрос открыт & Один опыт: во сне межклеточное пространство на $60\%$ шире и уборка быстрее. Другой, с иным способом измерения: во сне и под наркозом уборка заметно медленнее. Результаты противоречат друг другу & \cite{Xie2013,Miao2024} & гипотеза \\
20 & Местный сон: участки коры бодрствующего животного ненадолго проваливаются в тишину & У крыс после долгого бодрствования нейроны отдельного участка коры коротко замолкают, как в медленном сне, с медленной волной в местной ЭЭГ; в такие моменты крыса чаще ошибается в задаче & \cite{Vyazovskiy2011} & измерено \\
\bottomrule
\end{longtable}}

\section{Глава~\ref{sec:livingmachine}. Машина из живых нейронов}

Поведение культур на решётках электродов измерено в прямых опытах с записью и стимуляцией; механизм залпов опирается на модель истощения синапсов и опыты на микрокультурах, а утверждения о дофамине в чашке --- на единичные опыты, часть которых сами авторы считают лишь демонстрацией.

{\footnotesize
\setlength{\tabcolsep}{3pt}\renewcommand{\arraystretch}{1.15}
\begin{longtable}{>{\raggedright}p{0.5cm}>{\raggedright}p{4.0cm}>{\raggedright}p{5.6cm}>{\raggedright}p{2.3cm}>{\raggedright\arraybackslash}p{1.7cm}}
\toprule
№ & Утверждение & Опыт и что измерено & Источник & Статус \\
\midrule\endfirsthead
\toprule
№ & Утверждение & Опыт и что измерено & Источник & Статус \\
\midrule\endhead
1 & Культура на чипе: в основном \nt{ГЛУТ}, примерно пятая часть \nt{ГАМК} & Сеть из тысяч нейронов на решётке электродов --- стандартный опыт (строка~3). Для доли \nt{ГАМК} «около пятой части» глава ссылки не даёт, и проверенного источника для этого числа в культурах коры мы не нашли; доля зависит от препарата: в культурах гиппокампа \nt{ГАМК}-нейронов около $6\%$ & \cite{Benson1994} & измерено \\
2 & Органоиды: трёхмерная нервная ткань размером около миллиметра из стволовых клеток человека & Из стволовых клеток вырастили трёхмерные органоиды с несколькими областями мозга, в том числе корой со слоями предшественников и зрелыми нейронами & \cite{Lancaster2013} & измерено \\
3 & Без входа культура каждые несколько секунд вспыхивает залпом & $58$ культур коры плотностью от $3\,000$ до $50\,000$ нейронов записывали пять недель ($963$ получасовые записи): после двух недель почти во всех культурах преобладают залпы всей сети; интервалы между залпами --- от $1$ до $300$~с и сильно зависят от препарата & \cite{Wagenaar2006} & измерено \\
4 & Залп кончается истощением медиатора, интервал $T_{\text{залп}}\approx\tau_{\text{вос}}\ln\frac{1}{1-x^*}$~\eqref{eq:burstinterval} & Формула выведена в главе. Модель: сеть с истощающимися синапсами сама порождает залпы всей сети. Опыт на микрокультурах из $4$--$30$ нейронов: длительность залпа зависит от времени после предыдущего, истощение пузырьков медиатора отменяет залпы; вывод авторов --- залп ограничен запасом медиатора. Восстановление там --- десятки секунд, а не $0.8$~с & вывод в главе; \cite{Tsodyks2000,CohenSegal2011,Tsodyks1997} & теория \\
5 & «Учитель, который замолкает»: стимуляция выключается, когда появился нужный ответ, и ответ закрепляется & Культуру стимулировали с частотой $0.3$--$1$~Гц, пока через $50\pm10$~мс после стимула не появится выбранный ответ, затем стимуляцию выключали на $5$~мин; после нескольких циклов нужный ответ вызывался сразу; строили кривые обучения & \cite{Shahaf2001} & измерено \\
6 & Культура играет в теннис; серии отбитых мячей удлиняются только при обратной связи & Подробно --- в главе~\ref{sec:dishbrain} и приложении к ней & \cite{Kagan2022} & измерено \\
7 & Что культура учится предсказывать свой вход --- гипотеза; громкие слова о «разумности» оспорены & Принцип свободной энергии --- теоретическое предложение; прямого опыта, отличающего его от более простых объяснений, нет. Тридцать исследователей в ответной статье указали, что изменение поведения в петле не показывает ни разума, ни ощущений & \cite{Friston2010,Balci2023} & гипотеза \\
8 & Культура водит виртуальное тело к цели при подобранном рисунке стимулов & Программа сама подбирала обучающие рисунки стимуляции по текущему результату; за десятки минут сеть достигала заданных состояний, пластичность держалась $80$~мин после $2$~ч обучения. Те же стимулы, поданные без связи с поведением, не действовали & \cite{Bakkum2008} & измерено \\
9 & Резервуар: обучается только линейное считывание $y=W_{\text{вых}}\mathbf r$~\eqref{eq:reservoir} & Теория вычисления на резервуаре: любая нелинейная система с затухающей памятью плюс обученный линейный выход & \cite{Maass2002,JaegerHaas2004} & теория \\
10 & Органоид как резервуар различает говорящих с точностью около $78\%$ & Звуки речи нескольких говорящих подавали органоиду рисунками тока на решётке электродов; обученное считывание различало говорящего с точностью около $78\%$. Оговорка: авторы сообщают, что при тренировке менялась и функциональная связность органоида, и называют это обучением без учителя в самом органоиде, --- обучение не целиком снаружи & \cite{Cai2023} & измерено \\
11 & Миллион нейронов тратит на импульсы около $10^{-4}$~Вт~\eqref{eq:dishpower} & Оценка: цена импульса $10^{-10}$~Дж из энергетического бюджета коры~\eqref{eq:energy}, умноженная на число импульсов; прямого измерения мощности культуры нет & вывод в главе; \cite{Attwell2001} & теория \\
12 & Дофамин по вспышке света: $1$~мМ клеточного дофамина, $240$ повторов, число вызванных импульсов росло & На удалённой платформе с органоидами дофамин в светочувствительной клетке ($1$~мМ) освобождали ультрафиолетом ($365$~нм, $0.8$~с), когда стимул вызывал больше импульсов; за $240$ шагов (около $5$~ч) число импульсов в целом росло. Авторы пишут, что по одному опыту связь с дофамином установить нельзя; контролей нет & \cite{Jordan2024} & гипотеза \\
13 & Дофамин от живых клеток меняет вес органического транзистора надолго и обратимо & Клетки, выделяющие дофамин, выращены над органическим нейроморфным элементом; окисление дофамина у электрода меняет проводимость канала; показаны долговременное изменение веса и его возврат & \cite{Keene2020} & измерено \\
14 & Органоиды с работающими \nt{ДОФА}-нейронами существуют; как учитель их дофамин не использовали & В органоидах из стволовых клеток человека найдены электрически активные зрелые \nt{ДОФА}-нейроны и выработка дофамина. Опытов, где этот дофамин учит другую сеть, в литературе мы не нашли & \cite{Jo2016} & измерено \\
15 & Органоид балансирует шестом: учебные стимулы от программы лучше случайных, без закрепления, через \nt{ГЛУТ}-синапсы & Органоиды коры мыши в замкнутой петле с задачей «шест на тележке»: у большинства органоидов сигналы, выбранные обучением с подкреплением, дали лучший результат, чем случайные или никакие; после $45$~мин отдыха улучшение не сохранялось; блокада AMPA- и NMDA-рецепторов его отменяла & \cite{Robbins2026} & измерено \\
16 & Без регистров управления сеть на стенде не может сама решить, что считать успехом (утверждение~\ref{prop:livingmachine}) & Во всех опытах строк 5--15 учебный сигнал задаёт экспериментатор или программа; опыта, где культура сама выработала бы критерий успеха, нет --- это вывод из отсутствия, а не опыт & --- & гипотеза \\
\bottomrule
\end{longtable}}

\section{Глава~\ref{sec:dishbrain}. DishBrain}

Устройство стенда и результаты игры взяты из одной экспериментальной работы и её продолжения; формулы шума и дрейфа к хорошим настройкам выведены в главе и проверены моделью книги, а механизм обучения в чашке не установлен.

{\footnotesize
\setlength{\tabcolsep}{3pt}\renewcommand{\arraystretch}{1.15}
\begin{longtable}{>{\raggedright}p{0.5cm}>{\raggedright}p{4.0cm}>{\raggedright}p{5.6cm}>{\raggedright}p{2.3cm}>{\raggedright\arraybackslash}p{1.7cm}}
\toprule
№ & Утверждение & Опыт и что измерено & Источник & Статус \\
\midrule\endfirsthead
\toprule
№ & Утверждение & Опыт и что измерено & Источник & Статус \\
\midrule\endhead
1 & Стенд: около $800\,000$ клеток коры мыши или около $10^6$ клеток человека на чипе; $26\,000$ электродов на $8$~мм$^2$, запись до $1024$, стимуляция через $8$ & Методы работы: чип MaxOne, $26\,000$ платиновых электродов на $8$~мм$^2$, запись $1024$ каналов по $20\,000$ отсчётов в секунду; стимулировать технически можно до $32$ электродов, независимо удалось $8$; мышиные культуры использовали примерно с $10$~суток & \cite{Kagan2022} & измерено \\
2 & Петля с тактом $10$~мс: импульсы двигательных областей двигают ракетку (рис.~\ref{fig:dishloop}) & Импульсы считаются окнами по $10$~мс; задержка от импульса до ответного стимула около $5$~мс, её задаёт буфер аппаратуры & \cite{Kagan2022} & измерено \\
3 & Вход: код местом (какой из $8$ электродов) и частотой $4$--$40$~имп/с & В основном опыте частота стимуляции растёт линейно от $4$~Гц, когда мяч у дальней стенки, до $40$~Гц у ракетки; амплитуда импульсов мяча --- $75$~мВ; импульсы двухфазные прямоугольные & \cite{Kagan2022} & измерено \\
4 & Выход $\Delta p=c\,(n_\uparrow-n_\downarrow)$~\eqref{eq:dishreadout}, шум счёта за окно $1/\sqrt{RT}$ & Правило разности двух областей описано в работе (с весами областей по активности), точная формула не дана. Оценка шума --- следствие пуассоновского счёта, выведена в главе & \cite{Kagan2022}; вывод в главе & теория \\
5 & Учитель: после удара --- предсказуемый стимул $100$~имп/с, $0.1$~с; после промаха --- непредсказуемый, $150$~мВ, $5$~имп/с, $4$~с & Методы: после удара $75$~мВ, $100$~Гц, $100$~мс на все $8$ электродов сразу; после промаха $150$~мВ, $5$~Гц в случайные электроды в случайные моменты в течение $4$~с. Дофамина в культуре нет & \cite{Kagan2022} & измерено \\
6 & Сеть действует так, чтобы её вход стал предсказуемым & Принцип свободной энергии --- теоретическое предложение, из которого авторы выводили протокол; опыт с ним согласуется, но не отличает его от более простых механизмов (строки~7--8) & \cite{Friston2010,Kagan2022} & гипотеза \\
7 & Если неудача встряхивает сеть, а удача нет, время в настройке $\rho\propto1/P_{\text{пр}}$~\eqref{eq:dishdensity} (утверждение~\ref{prop:dishscramble}) & Следует из баланса потоков в марковской цепи с симметричными толчками, выведено в главе. Прямой проверки на культуре нет & вывод в главе & теория \\
8 & Модель «перемешивание при промахе» учится: доля ударов $0.21\to0.38$; при толчке после каждого розыгрыша $\approx0.12$, без толчков $0.19$ (рис.~\ref{fig:dishmodel}) & Расчёт, скрипт \texttt{models/dishbrain\_model.py}, $40$ моделей-культур по $2000$ розыгрышей: $0.21\to0.38$; $0.13\to0.11$; $0.19\to0.19$ & --- & модель \\
9 & Серии отбитых мячей удлиняются только у живых культур с обратной связью; контроли не учатся & $399$ сеансов по $20$~мин: $9$ культур мыши ($101$ сеанс), $11$ человека ($138$), $6$ чипов со средой ($80$), $20$ культур в покое ($42$), $3$ симуляции ($38$). Средняя длина розыгрыша за последние $15$~мин больше, чем за первые $5$, только у живых культур; у клеток, не являющихся нейронами (HEK293T), и у чипов со средой эффекта нет & \cite{Kagan2022} & измерено \\
10 & Без обратной связи и при тишине вместо стимула улучшения нет & $486$ сеансов за $3$ дня: «стимул» ($n=27$), «тишина» ($17$), «без обратной связи» ($15$). Рост во времени значим только в условии «стимул». Оговорка: в конце сеанса условие «тишина» всё же превосходило «без обратной связи», хоть и с меньшим эффектом & \cite{Kagan2022} & измерено \\
11 & Эффект невелик; учится ли сеть надолго --- открытый вопрос & Внутри сеанса улучшение надёжно, но, по словам самих авторов, обучение между сеансами за несколько дней устойчиво не наблюдалось & \cite{Kagan2022} & измерено \\
12 & Во время игры сеть приближается к критическому режиму, и чем ближе, тем лучше игра & Анализ лавин активности тех же культур: структурированный вход приближает сеть к критическому режиму, качество игры коррелирует с близостью к нему; но одной критичности без сведений о последствиях действий для обучения мало & \cite{Habibollahi2023} & измерено \\
13 & «Разумность» культуры не показана & Ответная статья тридцати исследователей: изменение поведения в замкнутой петле не доказывает ни разума, ни ощущений; опыта, который это показывал бы, нет & \cite{Balci2023} & гипотеза \\
14 & Предложенный четвёртый контроль: предсказуемый стимул после промаха той же длительности и энергии (раздел~\ref{sec:dishspec}) & Такой опыт не поставлен; это предложение книги & --- & гипотеза \\
\bottomrule
\end{longtable}}

